% Chapter 1, Section 3 _Linear Algebra_ Jim Hefferon
%  http://joshua.smcvt.edu/linearalgebra
%  2001-Jun-09
\section{简化阶梯形}
在掌握了高斯消元法的操作机制之后，
我们看到它可以用不止一种方式完成。
例如，从这个矩阵
\begin{equation*}
    \begin{mat}[r]
       2  &2  \\
       4  &3
    \end{mat}
\end{equation*}
出发，可以得到下面三个阶梯形矩阵中的任何一个。
\begin{equation*}
    \begin{mat}[r]
       2  &2  \\
       0  &-1
    \end{mat}
    \qquad
    \begin{mat}[r]
       1  &1  \\
       0  &-1
    \end{mat}
    \qquad
    \begin{mat}[r]
       2  &0  \\
       0  &-1
    \end{mat}
\end{equation*}
第一个由 $-2\rho_1+\rho_2$ 得到。
第二个来自先作 $(1/2)\rho_1$，再作 $-4\rho_1+\rho_2$。
第三个来自
$-2\rho_1+\rho_2$ 之后接 $2\rho_2+\rho_1$
（在第一个行倍加操作之后，矩阵已经成阶梯形
但这一步仍然是合法的行变换）。

在本章第一节中我们曾指出，
这带来了一些问题。
一个线性方程组的任意两个阶梯形版本，自由变量的
个数是否相同？
若是，
二者的自由变量是否恰好相同？
在本节中我们将
给出一种方法来判断一个线性方程组
能否由另一个经行变换得到。
这两个问题的答案都是“是”，
它们都将由此得出。








\subsection{高斯–若尔当消元}%
% Gaussian elimination coupled with back-substitution
% solves linear systems but it is not the only method possible.
下面是高斯消元法的一种推广，它有一些优点。

\begin{example} \label{exm:GJRedReadOffSol}
为求解
\begin{equation*}
  \begin{linsys}{3}
    x  &+  &y  &-  &2z  &=  &-2  \\
       &   &y  &+  &3z  &=  &7   \\
    x  &   &   &-  &z   &=  &-1
  \end{linsys}
\end{equation*}
我们可以像通常那样，先把它化简成阶梯形。
\begin{equation*}
  \grstep{-\rho_1+\rho_3}
    \begin{amat}[r]{3}
       1  &1  &-2 &-2  \\
       0  &1  &3  &7   \\
       0  &-1 &1  &1
    \end{amat}
  \grstep{\rho_2+\rho_3}
    \begin{amat}[r]{3}
       1  &1  &-2 &-2  \\
       0  &1  &3  &7   \\
       0  &0  &4  &8
    \end{amat}
\end{equation*}
我们可以继续进入第二阶段，
把首主元变成 \( 1 \)
\begin{equation*}
    \grstep{(1/4)\rho_3}
    \begin{amat}[r]{3}
       1  &1  &-2 &-2  \\
       0  &1  &3  &7   \\
       0  &0  &1  &2
    \end{amat}
\end{equation*}
然后再进入第三阶段，利用首主元
消去每列中所有其余的元素，
做法是向上倍加。
\begin{equation*}
  \grstep[2\rho_3+\rho_1]{-3\rho_3+\rho_2}
    \begin{amat}[r]{3}
       1  &1  &0  &2   \\
       0  &1  &0  &1   \\
       0  &0  &1  &2
    \end{amat}
  \grstep{-\rho_2+\rho_1}
    \begin{amat}[r]{3}
       1  &0  &0  &1   \\
       0  &1  &0  &1   \\
       0  &0  &1  &2
    \end{amat}
\end{equation*}
答案是 \( x=1 \)、\( y=1 \) 和 \( z=2 \)。
\end{example}
%<*Pivoting>
用一个元素去消去一列中其余的元素，称为关于该元素的
\definend{主元消去}\index{pivoting}。
%</Pivoting>

注意，第一阶段中的行倍加操作从左向右进行，
而第三阶段中的倍加操作则从右向左进行。

\begin{example}  \label{exm:GJRedReadOffSolTwo}
中间阶段把首主元变成 \( 1 \) 的那些操作
互不相干，所以我们可以把其中多个合并成单独的一步。
\begin{align*}
    \begin{amat}[r]{2}
       2   &1   &7   \\
       4   &-2  &6
    \end{amat}
  &\grstep{-2\rho_1+\rho_2}
  \begin{amat}[r]{2}
       2   &1   &7   \\
       0   &-4  &-8
    \end{amat}                                   \\
  &\grstep[(-1/4)\rho_2]{(1/2)\rho_1}
  \begin{amat}[r]{2}
       1   &1/2   &7/2   \\
       0   &1     &2
    \end{amat}                                    \\
  &\grstep{-(1/2)\rho_2+\rho_1}
  \begin{amat}[r]{2}
       1   &0   &5/2   \\
       0   &1   &2
    \end{amat}
\end{align*}
答案是 $x=5/2$ 和 $y=2$。
\end{example}

%<*GaussJordanReduction>
高斯消元法的这种推广称为
\definend{高斯–若尔当方法}\index{Gauss's Method!Gauss-Jordan Method}或
\definend{高斯–若尔当消元}。\index{linear equation!solution of!Gauss-Jordan}\index{Gauss's Method!Gauss-Jordan}
%</GaussJordanReduction>
% It goes past echelon form to a more refined, more specialized,
% matrix form.

\begin{definition}\label{def:RedEchForm}
%<*df:RedEchForm>
一个矩阵或线性方程组称为处于
\definend{简化阶梯形\/}\index{echelon form!reduced}\index{reduced echelon form}，
如果除了是阶梯形之外，每个首主元都是~$1$，
并且是其所在列中唯一的非零元素。
%</df:RedEchForm>
\end{definition}

\noindent
%<*CostRedEchForm>
用高斯–若尔当消元求解方程组的代价
是额外的算术运算。
好处是我们可以直接读出
解集的描述。
%</CostRedEchForm>

在任何阶梯形的方程组中，无论是简化阶梯形与否，
当方程组因存在矛盾方程而解集为空时，我们可以直接读出这一情况。
当方程组没有矛盾且每个变量都是某行的首变量时，
我们可以读出方程组的解集只有一个元素。
而且，当方程组没有矛盾且至少有一个变量是自由变量时，
我们可以读出方程组的解集有无穷多个元素。

然而，在简化阶梯形中，我们不仅能直接读出解集的
大小，还能读出它的描述。
当然，解集为空时描述它毫无困难。
\nearbyexample{exm:GJRedReadOffSol} 与~\ref{exm:GJRedReadOffSolTwo}
表明，在解集只有一个元素的情形，这个元素就在常数列中。
下一个例子演示如何读出无穷解集的参数化。

\begin{example}
\begin{multline*}
  \begin{amat}[r]{4}
     2  &6  &1  &2  &5  \\
     0  &3  &1  &4  &1  \\
     0  &3  &1  &2  &5
  \end{amat}
  \grstep{-\rho_2+\rho_3}
  \begin{amat}[r]{4}
     2  &6  &1  &2  &5  \\
     0  &3  &1  &4  &1  \\
     0  &0  &0  &-2 &4
  \end{amat}                                        \\
  \grstep[(1/3)\rho_2 \\ -(1/2)\rho_3]{(1/2)\rho_1}
  \repeatedgrstep[-\rho_3+\rho_1]{-(4/3)\rho_3+\rho_2}
  \repeatedgrstep{-3\rho_2+\rho_1}
  \begin{amat}[r]{4}
     1  &0  &-1/2  &0  &-9/2  \\
     0  &1  &1/3   &0  &3  \\
     0  &0  &0     &1  &-2
  \end{amat}
\end{multline*}
作为线性方程组，它就是
\begin{equation*}
  \begin{linsys}{4}
     x_1  &&     &-&1/2x_3  &&   &= &-9/2  \\
          &&x_2  &+&1/3x_3  &&   &= &3  \\
          &&     &&        &{}\hspace{.5em}{}&x_4 &= &-2
  \end{linsys}
\end{equation*}
因此解集的一个描述如下。
\begin{equation*}
  S=\set{\colvec{x_1 \\ x_2 \\ x_3 \\ x_4}
                        =\colvec[r]{-9/2 \\ 3 \\ 0 \\ -2}
                         +\colvec[r]{1/2 \\ -1/3 \\ 1 \\ 0}x_3
                        \suchthat x_3\in\Re}
\end{equation*}
\end{example}

可见，阶梯形并不是方程组的某种唯一最佳形式。
其他形式，比如简化阶梯形，各有其
优点和缺点。
与其把线性方程组（以及相应的矩阵）
想象成我们对其操作的对象，
总是朝着阶梯形这个目标推进，不如把它们
看成相互关联的，
即可以用行变换从一个得到另一个。
本小节余下的部分将展开这一想法。

\begin{lemma} \label{le:RowOpsRev}
%<*lm:RowOpsRev>
初等行变换都是可逆的。
%</lm:RowOpsRev>
\end{lemma}

\begin{proof}
%<*pf:RowOpsRev>
对任意矩阵 \( A \)，
对换两行的效果可以由再把它们对换回去而抵消，
用非零的 \( k \) 倍乘一行，可以由再用
$1/k$ 倍乘来撤销，
而把第 \( i \) 行的倍数加到第 \( j \) 行（$i\neq j$）
可以由从第 \( j \) 行减去第 \( i \) 行的同样倍数来撤销。
\begin{equation*}
      A
     \grstep{\rho_i\leftrightarrow\rho_j}
     \repeatedgrstep{\rho_j\leftrightarrow\rho_i}
      A
  \qquad
        A
     \grstep{k\rho_i}
     \repeatedgrstep{(1/k)\rho_i}
      A
  \qquad
        A
     \grstep{k\rho_i+\rho_j}
     \repeatedgrstep{-k\rho_i+\rho_j}
      A
\end{equation*}
%</pf:RowOpsRev>
（我们需要 $i\neq j$ 这一条件；
见 \nearbyexercise{exer:INotJMakesRowOpsRev}。）
\end{proof}

再说一次，我们正在发展、如今又得到引理支持的观点是：
“化简为”这一说法有误导性：~当
\( A\longrightarrow B \) 时，我们不应把 \( B \) 看作
在~\( A \) 之后，或比~$A$ 更简单。
相反，我们应当把这两个矩阵看作相互关联的。
下面是一幅图。
它把本节开头出现的那些矩阵以及它们的
简化阶梯形版本画成
一簇，彼此可互相化简。
\begin{center}
  \includegraphics{gr/mp/ch1.28}
\end{center}

%<*EquivMatrices>
我们说彼此可互相化简的矩阵，关于行可化简这一关系是等价的。
下一个结果将利用等价的定义
为此提供依据。\appendrefs{equivalence relations}
%</EquivMatrices>

\begin{lemma} \label{lm:ReducesToIsEqRel}
%<*lm:ReducesToIsEqRel>
在矩阵之间，“化简为”是一个等价关系。
%</lm:ReducesToIsEqRel>
\end{lemma}

\begin{proof}
%<*pf:ReducesToIsEqRel0>
我们必须验证以下条件：
(i)~自反性，即任何矩阵都化简到它自身；
(ii)~对称性，即若 \( A \) 化简为 \( B \)，
   则 \( B \) 化简为 \( A \)；
以及 (iii)~传递性，即若 \( A \) 化简为 \( B \) 且
      \( B \) 化简为 \( C \)，则 \( A \) 化简为 \( C \)。
%</pf:ReducesToIsEqRel0>

%<*pf:ReducesToIsEqRel1>
自反性是容易的；任何矩阵经零步操作即化简到它自身。

由前面的引理，这一关系是对称的\Dash 若
\( A \) 经某些行变换化简为 \( B \)，
那么 \( B \) 也可以通过将这些操作反过来而化简为 \( A \)。
%</pf:ReducesToIsEqRel1>

%<*pf:ReducesToIsEqRel2>
对于传递性，设 \( A \) 化简为 \( B \)，
且 \( B \) 化简为 \( C \)。
把 $A \rightarrow\cdots\rightarrow B$ 的化简步骤
接上 $B \rightarrow\cdots\rightarrow C$ 的化简步骤，
就得到从 \( A \) 到 \( C \) 的化简。
%</pf:ReducesToIsEqRel2>
\end{proof}

\begin{definition} \label{df:RowEquivalence}
%<*df:RowEquivalence>
经初等行变换而可互相化简的两个矩阵是\definend{行等价的}。\index{matrix!row equivalence}%
\index{row equivalence}\index{equivalence relation!row equivalence}
%</df:RowEquivalence>
\end{definition}

%<*RowEquivalanceClasses>
下面的图把全体矩阵的集合画成一个盒子。
在这个盒子中，每个矩阵位于一个类中。
两个矩阵在同一个类中当且仅当它们可互相化简。
这些类互不相交\Dash 没有矩阵属于两个不同的类。
我们就这样把矩阵的集合划分成了
\definend{行等价类}。\appendrefs{partitions and class representatives}\index{partition!row equivalence classes}
%</RowEquivalanceClasses>
\begin{center}
  \includegraphics{gr/mp/ch1.27}
\end{center}
\noindent 其中的一个类就是上面画出的、本节开头那些相互关联的矩阵组成的簇
（它包含所有非奇异的 $\nbyn{2}$ 矩阵）。

下一小节将证明一个矩阵的简化阶梯形是
唯一的。
用行等价关系来重新表述，
我们将证明每个矩阵都
与一个且仅一个简化阶梯形矩阵行等价。
用划分的语言来说，我们要证明的是：~每个
等价类恰含一个且仅一个简化阶梯形矩阵。
于是每个简化阶梯形矩阵都可以作为它
所在类的代表元。

\begin{exercises}
   \recommended \item 
     用高斯–若尔当消元求解每个方程组。
     \begin{exparts*}
        \partsitem \(
          \begin{linsys}[t]{2}
               x  &+  &y  &=  &2  \\
               x  &-  &y  &=  &0  
          \end{linsys}   \)
        \partsitem \(
          \begin{linsys}[t]{3}
               x  &   &   &-  &z  &=  &4  \\
              2x  &+  &2y &   &   &=  &1  
          \end{linsys}  \)
        \partsitem  \(
           \begin{linsys}[t]{2}
               3x  &-  &2y  &=  &1  \\
               6x  &+  &y   &=  &1/2 
           \end{linsys}  \)
        \partsitem \(
           \begin{linsys}[t]{3}
              2x  &-  &y  &  &  &= &-1  \\
               x  &+  &3y &- &z &= &5   \\
                  &   &y  &+ &2z&= &5   
           \end{linsys} \)
     \end{exparts*}
     \begin{answer}
       这些解答只给出了高斯–若尔当消元的过程。
       有了它，描述解集就很容易了。
       \begin{exparts}
         \partsitem 解集只含一个元素。
           \begin{multline*}
             \begin{amat}[r]{2}
               1  &1  &2  \\
               1  &-1 &0
             \end{amat}
             \grstep{-\rho_1+\rho_2}
             \begin{amat}[r]{2}
               1  &1  &2  \\
               0  &-2 &-2
             \end{amat}                          \\                       
             \grstep{-(1/2)\rho_2}
             \begin{amat}[r]{2}
               1  &1  &2  \\
               0  &1  &1
             \end{amat}                         
             \grstep{-\rho_2+\rho_1}
             \begin{amat}[r]{2}
               1  &0  &1  \\
               0  &1  &1
             \end{amat}
           \end{multline*}
         \partsitem 解集含有一个参数。
            \begin{equation*}
             \begin{amat}[r]{3}
               1  &0  &-1  &4  \\
               2  &2  &0   &1
             \end{amat}
             \grstep{-2\rho_1+\rho_2}
             \begin{amat}[r]{3}
               1  &0  &-1  &4  \\
               0  &2  &2   &-7
             \end{amat}                    
             \grstep{(1/2)\rho_2}
             \begin{amat}[r]{3}
               1  &0  &-1  &4  \\
               0  &1  &1   &-7/2
             \end{amat}
           \end{equation*}
         \partsitem 有唯一解。
           \begin{multline*}
             \begin{amat}[r]{2}
               3  &-2  &1  \\
               6  &1   &1/2
             \end{amat}
             \grstep{-2\rho_1+\rho_2}
             \begin{amat}[r]{2}
               3  &-2  &1  \\
               0  &5   &-3/2
             \end{amat}                            \\                       
             \grstep[(1/5)\rho_2]{(1/3)\rho_1}
             \begin{amat}[r]{2}
               1  &-2/3&1/3 \\
               0  &1   &-3/10
             \end{amat}                       
             \grstep{(2/3)\rho_2+\rho_1}
             \begin{amat}[r]{2}
               1  &0   &2/15 \\
               0  &1   &-3/10
             \end{amat}
          \end{multline*}
         \partsitem 第二步作一次对换两行的变换，会使算术更简单。
          \begin{multline*}
           \begin{amat}[r]{3}
             2  &-1  &0  &-1  \\
             1  &3   &-1 &5   \\
             0  &1   &2  &5
           \end{amat}
           \grstep{-(1/2)\rho_1+\rho_2}
           \begin{amat}[r]{3}
             2  &-1  &0  &-1   \\
             0  &7/2 &-1 &11/2 \\
             0  &1   &2  &5
           \end{amat}                                \\
           \begin{aligned}
           &\grstep{\rho_2\leftrightarrow\rho_3}
           \begin{amat}[r]{3}
             2  &-1  &0  &-1   \\
             0  &1   &2  &5    \\
             0  &7/2 &-1 &11/2
           \end{amat}                      
             \grstep{-(7/2)\rho_2+\rho_3}
             \begin{amat}[r]{3}
               2  &-1  &0  &-1   \\
               0  &1   &2  &5    \\
               0  &0   &-8 &-12
             \end{amat}                               \\     
             &\grstep[-(1/8)\rho_2]{(1/2)\rho_1}
             \begin{amat}[r]{3}
               1  &-1/2&0  &-1/2 \\
               0  &1   &2  &5    \\
               0  &0   &1  &3/2
             \end{amat}                                
             \grstep{-2\rho_3+\rho_2}
             \begin{amat}[r]{3}
               1  &-1/2&0  &-1/2 \\
               0  &1   &0  &2    \\
               0  &0   &1  &3/2
             \end{amat}                                    \\             
             &\grstep{(1/2)\rho_2+\rho_1}
             \begin{amat}[r]{3}
               1  &0   &0  &1/2  \\
               0  &1   &0  &2    \\
               0  &0   &1  &3/2
             \end{amat}
           \end{aligned}
         \end{multline*}
       \end{exparts}  
     \end{answer}
   \item 做高斯–若尔当消元。
     \begin{exparts*}
       \partsitem
         $\begin{linsys}[t]{3}
           x  &+ &y &- &z &= &3 \\
           2x &- &y  &- &z &= &1 \\
           3x &+  &y  &+ &2z &= &0
         \end{linsys}$
       \partsitem
         $\begin{linsys}[t]{3}
           x  &+ &y  &+ &2z  &= &0 \\
           2x  &- &y  &+  &z &= &1 \\
           4x &+ &y  &+ &5z &= &1  
         \end{linsys}$
     \end{exparts*}
     \begin{answer}
       \begin{exparts}
         \partsitem
           \begin{multline*}
             \begin{amat}{3}
               1 &1  &-1 &3 \\
               2 &-1 &-1 &1 \\
               3 &1  &2  &0
             \end{amat}
             \grstep[-3\rho_1+\rho_3]{-2\rho_1+\rho_2}
             \begin{amat}{3}
               1 &1  &-1 &3 \\
               0 &-3 &1  &-5 \\
               0 &-2  &5  &-9
             \end{amat}                             \\
           \begin{aligned}   
             &\grstep{-(2/3)\rho_2+\rho_3}
             \begin{amat}{3}
               1 &1  &-1    &3 \\
               0 &-3 &1     &-5 \\
               0 &0  &13/3  &-17/3
             \end{amat}                                 
             \grstep[(3/13)\rho_3]{-(1/3)\rho_2}
             \begin{amat}{3}
               1 &1  &-1    &3 \\
               0 &1  &-1/3    &5/3 \\
               0 &0  &1      &-17/13
             \end{amat}                                \\
             &\grstep[(1/3)\rho_3+\rho_2]{\rho_3+\rho_1}
             \begin{amat}{3}
               1 &1  &0    &22/13 \\
               0 &1  &0    &16/13 \\
               0 &0  &1      &-17/13
             \end{amat}
             \grstep{-\rho_2+\rho_1}
             \begin{amat}{3}
               1 &0  &0    &6/13 \\
               0 &1  &0    &16/13 \\
               0 &0  &1      &-17/13
             \end{amat}
           \end{aligned}
           \end{multline*}
         \partsitem
           \begin{multline*}
             \begin{amat}{3}
               1 &1  &2 &0 \\
               2 &-1 &1 &1 \\
               4 &1  &5 &1
             \end{amat}
             \grstep[-4\rho_1+\rho_3]{-2\rho_1+\rho_2}
             \begin{amat}{3}
               1 &1  &2  &0 \\
               0 &-3 &-3 &1 \\
               0 &-3 &-3 &1
             \end{amat}
             \grstep{-\rho_2+\rho_3}
             \begin{amat}{3}
               1 &1  &2  &0 \\
               0 &-3 &-3 &1 \\
               0 &0  &0  &0
             \end{amat}                              \\
             \grstep{-(1/3)\rho_2}
             \begin{amat}{3}
               1 &1  &2  &0 \\
               0 &1  &1 &-1/3 \\
               0 &0  &0  &0
             \end{amat}
             \grstep{-\rho_2+\rho_1}
             \begin{amat}{3}
               1 &0  &1  &1/3 \\
               0 &1  &1 &-1/3 \\
               0 &0  &0  &0
             \end{amat}
           \end{multline*}
       \end{exparts}
     \end{answer}
   \recommended \item 
     求出每个矩阵的简化阶梯形。
     \begin{exparts*}
       \partsitem \( \begin{mat}[r]
           2  &1  \\
           1  &3
         \end{mat}  \)
       \partsitem \( \begin{mat}[r]
           1  &3  &1  \\
           2  &0  &4  \\
          -1  &-3 &-3
         \end{mat}  \)
       \partsitem \( \begin{mat}[r]
           1  &0  &3  &1  &2  \\
           1  &4  &2  &1  &5  \\
           3  &4  &8  &1  &2
         \end{mat}  \)
       \partsitem \( \begin{mat}[r]
           0  &1  &3  &2  \\
           0  &0  &5  &6  \\
          1  &5  &1  &5
        \end{mat}  \)
    \end{exparts*}
    \begin{answer}
      使用高斯–若尔当消元。
      \begin{exparts}
        \partsitem 简化阶梯形除了由1构成的对角线以外全为0。
          \begin{equation*}
            \grstep{-(1/2)\rho_1+\rho_2}
            \begin{mat}[r]
              2  &1  \\
              0  &5/2
            \end{mat}
            \grstep[(2/5)\rho_2]{(1/2)\rho_1}
            \begin{mat}[r]
              1  &1/2\\
              0  &1
            \end{mat}                      
            \grstep{-(1/2)\rho_2+\rho_1}
            \begin{mat}[r]
              1  &0  \\
              0  &1
            \end{mat}
          \end{equation*}
        \partsitem 与上一题一样，简化阶梯形
          除了由1构成的对角线以外全是0。
          \begin{multline*}
            \grstep[\rho_1+\rho_3]{-2\rho_1+\rho_2}
            \begin{mat}[r]
              1  &3  &1  \\
              0  &-6 &2  \\
              0  &0  &-2
            \end{mat}
            \grstep[-(1/2)\rho_3]{-(1/6)\rho_2}
            \begin{mat}[r]
              1  &3  &1     \\
              0  &1  &-1/3  \\
              0  &0  &1
            \end{mat}                                      \\              
            \grstep[-\rho_3+\rho_1]{(1/3)\rho_3+\rho_2}
            \begin{mat}[r]
              1  &3  &0     \\
              0  &1  &0     \\
              0  &0  &1
            \end{mat}                  
            \grstep{-3\rho_2+\rho_1}
            \begin{mat}[r]
              1  &0  &0     \\
              0  &1  &0     \\
              0  &0  &1
            \end{mat}
           \end{multline*}
        \partsitem 列数比行数多，所以我们得到的
          不只是由1构成的一条对角线。
          \begin{multline*}
            \grstep[-3\rho_1+\rho_3]{-\rho_1+\rho_2}
            \begin{mat}[r]
              1  &0  &3  &1  &2  \\
              0  &4  &-1 &0  &3  \\
              0  &4  &-1 &-2 &-4
            \end{mat}
            \grstep{-\rho_2+\rho_3}
            \begin{mat}[r]
              1  &0  &3  &1  &2  \\
              0  &4  &-1 &0  &3  \\
              0  &0  &0  &-2 &-7
            \end{mat}                           \\
            \grstep[-(1/2)\rho_3]{(1/4)\rho_2}
            \begin{mat}[r]
              1  &0  &3    &1  &2  \\
              0  &1  &-1/4 &0  &3/4  \\
              0  &0  &0    &1  &7/2
            \end{mat}                           
            \grstep{-\rho_3+\rho_1}
            \begin{mat}[r]
              1  &0  &3    &0  &-3/2  \\
              0  &1  &-1/4 &0  &3/4     \\
              0  &0  &0    &1  &7/2
            \end{mat}
          \end{multline*}
        \partsitem 与上一小题一样，这不是一个方阵。 
          \begin{multline*}
            \grstep{\rho_1\leftrightarrow\rho_3}
            \begin{mat}[r]
              1  &5  &1  &5  \\
              0  &0  &5  &6  \\
              0  &1  &3  &2
            \end{mat}
            \grstep{\rho_2\leftrightarrow\rho_3}
            \begin{mat}[r]
              1  &5  &1  &5  \\
              0  &1  &3  &2  \\
              0  &0  &5  &6
            \end{mat}                    \\
            \begin{aligned}
              &\grstep{(1/5)\rho_3}
              \begin{mat}[r]
                1  &5  &1  &5  \\
                0  &1  &3  &2  \\
                0  &0  &1  &6/5
              \end{mat}                  
              \grstep[-\rho_3+\rho_1]{-3\rho_3+\rho_2}
              \begin{mat}[r]
                1  &5  &0  &19/5  \\
                0  &1  &0  &-8/5  \\
                0  &0  &1  &6/5
              \end{mat}                   \\
              &\grstep{-5\rho_2+\rho_1}
              \begin{mat}[r]
                1  &0  &0  &59/5  \\
                0  &1  &0  &-8/5  \\
                0  &0  &1  &6/5
              \end{mat}
            \end{aligned}
          \end{multline*}
      \end{exparts}  
    \end{answer}
  \item 求出每一个的简化阶梯形。
    \begin{exparts*}
      \partsitem $
        \begin{mat}
          0  &2 &1 \\
          2  &-1 &1 \\
          -2 &-1 &0
        \end{mat}$
      \partsitem $
        \begin{mat}
          1 &3 &1 \\
          2 &6 &2 \\
         -1 &0 &0
        \end{mat}$
    \end{exparts*}
    \begin{answer}
      \begin{exparts}
        \partsitem 先作对换。
          \begin{equation*}
            \grstep{\rho_1\leftrightarrow\rho_2}
            \grstep{\rho_1+\rho_3}
            \grstep{\rho_2+\rho_3}
            \grstep[(1/2)\rho_2 \\ (1/2)\rho_3]{(1/2)\rho_1}
            \grstep[(-1/2)\rho_3+\rho_2]{(-1/2)\rho_3+\rho_1}
            \grstep{(1/2)\rho_2+\rho_1}
            \begin{mat}
              1 &0 &0 \\
              0 &1 &0 \\
              0 &0 &1
            \end{mat}
          \end{equation*}
        \partsitem 这里对换在中间进行。
        \begin{equation*}
          \grstep[\rho_1+\rho_3]{-2\rho_1+\rho_2}
          \grstep{\rho_2\leftrightarrow\rho_3}
          \grstep{(1/3)\rho_2}
          \grstep{-3\rho_2+\rho_1}
          \begin{mat}
            1 &0 &0   \\
            0 &1 &1/3 \\
            0 &0 &0
          \end{mat}
        \end{equation*}
      \end{exparts}
    \end{answer}
  \recommended \item 
    用高斯–若尔当消元求出每个解集，
    然后读出参数化。
    \begin{exparts}
      \partsitem \( \begin{linsys}[t]{3}
                  2x  &+  &y  &-  &z  &=  &1  \\
                  4x  &-  &y  &   &   &=  &3  
                  \end{linsys}  \)
      \partsitem \( \begin{linsys}[t]{4}
                   x  &   &   &-  &z  &   &   &=  &1  \\
                      &   &y  &+  &2z &-  &w  &=  &3  \\
                   x  &+  &2y &+  &3z &-  &w  &=  &7  
                    \end{linsys}  \)
      \partsitem \( \begin{linsys}[t]{4}
                   x  &-  &y  &+  &z  &   &   &=  &0  \\
                      &   &y  &   &   &+  &w  &=  &0  \\
                  3x  &-  &2y &+  &3z &+  &w  &=  &0  \\
                      &   &-y &   &   &-  &w  &=  &0  
                  \end{linsys}  \)
      \partsitem \( \begin{linsys}[t]{5}
                   a  &+  &2b &+  &3c &+  &d  &-  &e  &=  &1  \\
                  3a  &-  &b  &+  &c  &+  &d  &+  &e  &=  &3  
                  \end{linsys}  \)
    \end{exparts}
    \begin{answer}
      至于高斯消元的那一半，
      参见第一章第~I.2 节习题
      \nearbyexercise{exer:SlvMatNot} 的解答。
      \begin{exparts}
      \partsitem “若尔当”那一半如下进行。
        \begin{equation*}
          \grstep[-(1/3)\rho_2]{(1/2)\rho_1}
          \begin{amat}[r]{3}
            1  &1/2 &-1/2 &1/2  \\
            0  &1   &-2/3 &-1/3
          \end{amat}                           
          \grstep{-(1/2)\rho_2+\rho_1}
          \begin{amat}[r]{3}
            1  &0   &-1/6 &2/3  \\
            0  &1   &-2/3 &-1/3
          \end{amat}
        \end{equation*}
        解集为下面这个集合
        \begin{equation*}
          \set{\colvec[r]{2/3 \\ -1/3 \\ 0}
               +\colvec[r]{1/6 \\ 2/3 \\ 1}z
              \suchthat z\in\Re}
        \end{equation*}
      \partsitem 后一半是
        \begin{equation*}
          \grstep{\rho_3+\rho_2}
          \begin{amat}[r]{4}
            1  &0  &-1  &0  &1 \\
            0  &1  &2   &0  &3 \\
            0  &0  &0   &1  &0
          \end{amat}
        \end{equation*}
        于是解如下。
        \begin{equation*}
          \set{\colvec[r]{1 \\ 3 \\ 0 \\ 0}
               +\colvec[r]{1 \\ -2 \\ 1 \\ 0}z
              \suchthat z\in\Re}
        \end{equation*}
      \partsitem 这个若尔当一半
        \begin{equation*}
          \grstep{\rho_2+\rho_1}
          \begin{amat}[r]{4}
            1  &0  &1   &1  &0 \\
            0  &1  &0   &1  &0 \\
            0  &0  &0   &0  &0 \\
            0  &0  &0   &0  &0
          \end{amat}
        \end{equation*}
        给出
        \begin{equation*}
          \set{\colvec[r]{0 \\ 0 \\ 0 \\ 0}
               +\colvec[r]{-1 \\ 0 \\ 1 \\ 0}z
               +\colvec[r]{-1 \\ -1 \\ 0 \\ 1}w
              \suchthat z,w\in\Re}
        \end{equation*}
        （当然，我们可以在描述中省去零向量）。
      \partsitem “若尔当”那一半
        \begin{align*}
          &\grstep{-(1/7)\rho_2}
          \begin{amat}[r]{5}
            1  &2  &3   &1   &-1   &1  \\
            0  &1  &8/7 &2/7 &-4/7 &0
          \end{amat}                   \\
          &\grstep{-2\rho_2+\rho_1}
          \begin{amat}[r]{5}
            1  &0  &5/7 &3/7 &1/7  &1  \\
            0  &1  &8/7 &2/7 &-4/7 &0
          \end{amat}
        \end{align*}
        以这个解集收尾。
        \begin{equation*}
          \set{\colvec[r]{1 \\ 0 \\ 0 \\ 0 \\ 0}
               +\colvec[r]{-5/7 \\ -8/7 \\ 1 \\ 0 \\ 0}c
               +\colvec[r]{-3/7 \\ -2/7 \\ 0 \\ 1 \\ 0}d
               +\colvec[r]{-1/7 \\ 4/7 \\ 0 \\ 0 \\ 1}e
              \suchthat c,d,e\in\Re}
        \end{equation*}
    \end{exparts}
   \end{answer}
  \item 
    给出这个矩阵的两个不同的阶梯形版本。
    \begin{equation*}
      \begin{mat}[r]
        2  &1  &1  &3  \\
        6  &4  &1  &2  \\
        1  &5  &1  &5
      \end{mat}
    \end{equation*}
    \begin{answer}
      按常规作高斯消元法便得到一个：
      \begin{equation*}
        \grstep[-(1/2)\rho_1+\rho_3]{-3\rho_1+\rho_2}
        \begin{mat}[r]
          2  &1  &1  &3  \\
          0  &1  &-2 &-7 \\
          0  &9/2&1/2&7/2
        \end{mat}
        \grstep{-(9/2)\rho_2+\rho_3}
        \begin{mat}[r]
          2  &1  &1    &3  \\
          0  &1  &-2   &-7 \\
          0  &0  &19/2 &35
        \end{mat}
      \end{equation*}
      而任何外观上的改动，比如把最下面一行乘以 \( 2 \)，
      \begin{equation*}
        \begin{mat}[r]
          2  &1  &1    &3  \\
          0  &1  &-2   &-7 \\
          0  &0  &19   &70
        \end{mat}
      \end{equation*}
      便给出另一个。  
    \end{answer}
  \recommended \item \label{exer:PossRedEchFrms} 
    列出各种尺寸可能出现的简化阶梯形。
    \begin{exparts*}
      \partsitem \( \nbyn{2} \)
      \partsitem \( \nbym{2}{3} \)
      \partsitem \( \nbym{3}{2} \)
      \partsitem \( \nbyn{3} \)
    \end{exparts*}
    \begin{answer}
      在下面列出的各情形中，我们取 $a,b\in\Re$。
      因此，下面列出的某些标准形
      实际上包含了无穷多种情形。
      特别地，它们包括了 $a=0$ 与 $b=0$ 的情形。
      \begin{exparts}
        \partsitem 
          $\begin{mat}[r]
            0  &0  \\
            0  &0
          \end{mat}$,
          $\begin{mat}[r]
            1  &a  \\
            0  &0
          \end{mat}$, 
          $\begin{mat}[r]
            0  &1  \\
            0  &0
          \end{mat}$, 
          $\begin{mat}[r]
            1  &0  \\
            0  &1
          \end{mat}$
        \partsitem
          $\begin{mat}[r]
               0  &0  &0  \\
               0  &0  &0
             \end{mat}$,
          $\begin{mat}[r]
               1  &a  &b  \\
               0  &0  &0
             \end{mat}$,
          $\begin{mat}[r]
               0  &1  &a  \\
               0  &0  &0
             \end{mat}$,
          $\begin{mat}[r]
               0  &0  &1  \\
               0  &0  &0
             \end{mat}$,
          $\begin{mat}[r]
               1  &0  &a  \\
               0  &1  &b
             \end{mat}$,
          $\begin{mat}[r]
               1  &a  &0  \\
               0  &0  &1
             \end{mat}$,
          $\begin{mat}[r]
               0  &1  &0  \\
               0  &0  &1
             \end{mat}$
        \partsitem
          $\begin{mat}[r]
               0  &0  \\
               0  &0  \\
               0  &0
             \end{mat}$,
          $\begin{mat}[r]
               1  &a  \\
               0  &0  \\
               0  &0
             \end{mat}$,
          $\begin{mat}[r]
               0  &1  \\
               0  &0  \\
               0  &0
             \end{mat}$,
          $\begin{mat}[r]
               1  &0  \\
               0  &1  \\
               0  &0
             \end{mat}$
        \partsitem
          $\begin{mat}[r]
               0  &0  &0  \\
               0  &0  &0  \\
               0  &0  &0
             \end{mat}$,
          $\begin{mat}[r]
               1  &a  &b  \\
               0  &0  &0  \\
               0  &0  &0
             \end{mat}$,
          $\begin{mat}[r]
               0  &1  &a  \\
               0  &0  &0  \\
               0  &0  &0
             \end{mat}$,
          $\begin{mat}[r]
               0  &1  &0  \\
               0  &0  &1  \\
               0  &0  &0
             \end{mat}$,
          $\begin{mat}[r]
               0  &0  &1  \\
               0  &0  &0  \\
               0  &0  &0
             \end{mat}$,
          $\begin{mat}[r]
               1  &0  &a  \\
               0  &1  &b  \\
               0  &0  &0
             \end{mat}$,
          $\begin{mat}[r]
               1  &a  &0  \\
               0  &0  &1  \\
               0  &0  &0
             \end{mat}$,
          $\begin{mat}[r]
               1  &0  &0  \\
               0  &1  &0  \\
               0  &0  &1
             \end{mat}$
      \end{exparts}  
    \end{answer}
  \recommended \item  
    对非奇异矩阵作高斯–若尔当消元
    会得到什么结果？
    \begin{answer}
      非奇异的齐次线性方程组有唯一解。
      所以非奇异矩阵必定能化简成一个（方） 
      矩阵，它除了从左上到右下的对角线上是 \( 1 \) 之外
      全为 \( 0 \)，
      比如下面这些。
      \begin{equation*}
         \begin{mat}[r]
           1  &0  \\
           0  &1  \\
         \end{mat}
         \qquad
         \begin{mat}[r]
           1  &0  &0  \\
           0  &1  &0  \\
           0  &0  &1
         \end{mat}
      \end{equation*}  
    \end{answer}
 \item 判断下述每个关系是否是 $\nbyn{2}$ 矩阵
   集合上的等价关系。
   \begin{exparts}
     \partsitem 若两个矩阵第一行第一列的元素相同，
       则称这两个矩阵相关。
     % \partsitem two matrices are related if their four entries sum to the 
     %   same total
     \partsitem 若两个矩阵第一行第一列的元素相同，
       或者第二行第二列的元素相同，
       则称这两个矩阵相关。
   \end{exparts}
   \begin{answer}
     \begin{exparts}
       \partsitem  这是一个等价关系。
       
          若两个矩阵相关，我们可以写成 $M_1\sim M_2$。
          $\sim$~关系是自反的，因为任何矩阵与它自身
          有相同的 $1,1$ 元素。
          该关系是对称的，因为如果 $M_1$ 与~$M_2$ 有相同的 $1,1$
          元素，那么显然 $M_2$ 与~$M_1$ 也有相同的
          $1,1$ 元素。
          最后，该关系是传递的，因为如果 $M_1\sim M_2$，即两者
          有相同的 $1,1$ 元素，
          且 $M_2\sim M_3$，即两者
          也有相同的元素，那么三者都有相同的~$1,1$ 元素，
          从而 $M_1\sim M_3$。 
       % \partsitem  This is an equivalence.
       %    Write $M_1\sim M_2$ if they are related.

       %    This relation is reflexive because any matrix has the 
       %    same sum of entries as itself.
       %    The $\sim$~relation is symmetric because if $M_1$ has the same
       %    sum of 
       %    entries as~$M_2$, then also $M_2$ has the same sum of entries
       %    as~$M_1$.
       %    Finally, the relation is transitive because if $M_1\sim M_2$ so their
       %    entry sum is the same, 
       %    and $M_2\sim M_3$ so they have the 
       %    same entry sum as each other, then all three have the same entry sum, 
       %    and $M_1\sim M_3$. 
       \partsitem
          这不是等价关系，因为它不是传递的。
          下面第一个矩阵与第二个矩阵因 $1,1$ 元素而相关，
          第二个与第三个因 $2,2$~元素而相关。
          但第一个与第三个不相关。
          \begin{equation*}
            \begin{mat}
              1 &0 \\
              0 &0
            \end{mat}
            \quad
            \begin{mat}
              1 &0 \\
              0 &-1
            \end{mat}
            \quad
            \begin{mat}
              0 &0 \\
              0 &-1
            \end{mat}
          \end{equation*}
     \end{exparts}
   \end{answer}
 \item \cite{Cleary}
    考虑 $\nbyn{2}$ 矩阵集合上的如下关系：  
    如果 $A$ 中所有元素之和与 $B$ 中所有元素之和相同，
    我们就说 $A$ 与 $B$ \textit{总和相像}。  
    例如，零矩阵与第一行有两个七、
    第二行有两个负七的矩阵
    总和相像。
    证明或否定这是 $\nbyn{2}$ 矩阵集合
    上的等价关系。
    \begin{answer}
      它是等价关系。
      为证明这一点，我们必须验证该关系
      是自反的、对称的和传递的。

      设所有矩阵都是 $\nbyn{2}$ 的。
      自反性：我们注意到矩阵与它自身
      的元素之和相同。
      对称性：设 $A$ 与~$B$ 的元素之和相同，
      那么显然 $B$ 与~$A$ 的元素之和也相同。
      传递性同样不难\Dash 如果 $A$ 与 $B$ 的元素之和相同，
      且 $B$ 与 $C$ 的元素之和相同，
      那么 $A$ 与 $C$ 的元素之和相同。
    \end{answer}
 \item \label{exer:INotJMakesRowOpsRev}
   \nearbylemma{le:RowOpsRev} 的证明中提到了
   倍加操作上的 $i\neq j$ 条件。
   \begin{exparts}
     \partsitem 写出一个元素都非零的 $\nbyn{2}$ 矩阵，
        并证明 $-1\cdot\rho_1+\rho_1$ 操作
        不会被 $1\cdot\rho_1+\rho_1$ 抵消。
     \partsitem 扩充该引理的证明，明确指出
        它究竟在何处用到倍加的 $i\neq j$ 条件。
   \end{exparts}
   \begin{answer}
    \begin{exparts}
      \partsitem 例如，
        \begin{equation*}
          \begin{mat}[r]
            1  &2  \\
            3  &4  
          \end{mat}
          \grstep{-\rho_1+\rho_1}
          \begin{mat}[r]
            0  &0  \\
            3  &4  
          \end{mat}
          \grstep{\rho_1+\rho_1}
          \begin{mat}[r]
            0  &0  \\
            3  &4  
          \end{mat}
        \end{equation*}
        矩阵被改变了。
      \partsitem 这个操作
        \begin{equation*}
          \begin{mat}
            \vdotswithin{a_{i,1}}                     \\
            a_{i,1}  &\cdots  &a_{i,n}  \\
            \vdotswithin{a_{i,1}}                     \\
            a_{j,1}  &\cdots  &a_{j,n}  \\
            \vdotswithin{a_{i,1}}                     
          \end{mat}
          \grstep{k\rho_i+\rho_j}
          \begin{mat}
            \vdotswithin{a_{i,1}}                                      \\
            a_{i,1}           &\cdots  &a_{i,n}          \\
            \vdotswithin{a_{i,1}}                                      \\
            ka_{i,1}+a_{j,1}  &\cdots  &ka_{i,n}+a_{j,n}  \\
            \vdotswithin{a_{i,1}}                     
          \end{mat}
        \end{equation*}      
        由于 $i\neq j$ 的限制而使第$i$行保持不变。
        因为第$i$行未变，所以这个操作
        \begin{equation*}                                        
          \grstep{-k\rho_i+\rho_j}
          \begin{mat}
            \vdotswithin{a_{i,1}}                                      \\
            a_{i,1}           &\cdots  &a_{i,n}          \\
            \vdotswithin{a_{i,1}}                                      \\
            -ka_{i,1}+ka_{i,1}+a_{j,1}  &\cdots &-ka_{i,n}+ka_{i,n}+a_{j,n} \\
            \vdotswithin{a_{i,1}}                     
          \end{mat}
        \end{equation*}
        把第$j$行还原到它原来的状态。
        % (If $i=j$ then the third matrix would have entries of the 
        % form $-k(ka_{i,j}+a_{i,j})+ka_{i,j}+a_{i,j}$.)
    \end{exparts}
   \end{answer}
 \recommended \item \cite{Cleary}
  考虑某个班级学生的集合。  
  下列关系中哪些是等价关系？  
  对每个答案至少用一句话解释。
  \begin{exparts}
    \item  两个学生 $x, y$ 相关，
      如果 $x$ 修过的数学课
      至少与 $y$ 一样多。
    \item 学生 $x, y$ 相关，如果他们的名字
      以同一个字母开头。
  \end{exparts}
  \begin{answer}
    要成为等价关系，每个关系都必须是自反的、对称的和
    传递的。
    \begin{exparts}
      \item 这个关系
        不是对称的，因为如果 $x$ 修了 $4$~门课而 $y$
        修了 $3$ 门，那么 $x$ 与 $y$ 相关，但 $y$ 与 $x$
        不相关。
      \item 这是自反的，因为 $x$ 的名字与 $x$ 自己的名字
        以同一个字母开头。
        它是对称的，因为如果 $x$ 的名字与 $y$ 的名字
        以同一个字母开头，那么 $y$ 的名字也与~$x$ 的以同一个字母开头。
        它还是传递的，因为如果 $x$ 的名字与~$y$ 的以同一个字母开头，
        且 $y$ 的名字与 $z$ 的
        以同一个字母开头，那么 $x$ 的名字与 $z$ 的以同一个字母开头。
        所以它是等价关系。
    \end{exparts}
  \end{answer}
  \item
  证明下述每一个都是 $\nbyn{2}$ 矩阵集合
  上的等价关系。
  描述各个等价类。
  \begin{exparts}
    \item 如果两个矩阵沿对角线的乘积相同，
      也就是说，左上角元素与
      右下角元素的乘积相等，则这两个矩阵相关。
    \item 如果两个矩阵都至少有一个元素
      为~$1$，
      或者两者都没有，则这两个矩阵相关。
  \end{exparts}
  \begin{answer}
    对每一个我们都必须验证自反性、对称性和
    传递性这三个条件。
    \begin{exparts}
      \item 显然，任何矩阵沿对角线的乘积都与它自身相同，
        所以该关系是自反的。
        该关系是对称的，因为如果 $A$ 沿对角线的乘积
        与~$B$ 的相同，
        即 $a_{1,1}\cdot a_{2,2}=b_{1,1}\cdot b_{2,2}$，
        那么 $B$ 的乘积也与~$A$ 的相同。
 
        传递性类似：设 $A$ 的乘积为~$r$，
        且等于 $B$ 的乘积。
        又设 $B$ 的乘积等于 $C$ 的乘积。
        那么三者的乘积都是~$r$，即 $A$ 的乘积等于~$C$ 的乘积。

        每个实数对应一个等价类，也就是说，
        该等价类包含所有沿对角线
        乘积为这个实数的 $\nbyn{2}$ 矩阵。
      \item 自反性：如果矩阵 $A$ 有一个元素为~$1$，那么它与
        自身相关；如果没有，那么它同样与
        自身相关。
        对称性也分两种情形：设 $A$ 与~$B$ 相关。
        如果 $A$ 有一个元素为~$1$，那么~$B$ 也有，于是 $B$ 与~$A$ 相关。
        如果 $A$ 没有~$1$，那么 $B$ 也没有，同样 $B$ 与
        $A$ 相关。
 
        传递性：设 $A$ 与~$B$ 相关，$B$ 与~$C$ 相关。
        如果 $A$ 有一个元素为~$1$，那么~$B$ 也有；又因 $B$ 与~$C$ 相关，
        所以~$C$ 也有，
        从而 $A$ 与~$C$ 相关。
        同样，如果 $A$ 没有~$1$，那么~$B$ 也没有，因而
        ~$C$ 也没有，结论同样是 $A$ 与~$C$ 相关。 

        恰有两个等价类：一个包含任何
        至少有一个元素为~$1$ 的 $\nbyn{2}$ 矩阵，
        另一个包含所有不含 $1$ 的矩阵。
    \end{exparts}
  \end{answer}
  \item 证明下述每一个都不是 $\nbyn{2}$ 矩阵
    集合上的等价关系。
    \begin{exparts}
      \item 两个矩阵 $A,B$ 相关，如果 $a_{1,1}=-b_{1,1}$。
      \item 如果两个矩阵的元素之和相差在 $5$ 以内，
        即当
        $\absval{(a_{1,1}+\cdots+a_{2,2})-(b_{1,1}+\cdots+b_{2,2})}<5$ 时 $A$ 与~$B$ 相关，则称这两个矩阵相关。
    \end{exparts}
    \begin{answer}
      \begin{exparts}
        \item 这个关系不是自反的。
          例如，任何左上角元素为~$1$ 的矩阵
          都与自身不相关。
        \item 这个关系不是传递的。
          对下面这三个矩阵，$A$ 与~$B$ 相关，$B$ 与~$C$ 相关，
          但 $A$ 与~$C$ 不相关。
          \begin{equation*}
            A=\begin{mat}
              0 &0 \\
              0 &0
            \end{mat},\quad
            B=\begin{mat}
              4 &0 \\
              0 &0
            \end{mat},\quad
            C=\begin{mat}
              8 &0 \\
              0 &0
            \end{mat},\quad
          \end{equation*}
      \end{exparts}
    \end{answer}
\end{exercises}




















\subsection{线性组合引理}
我们将以证明每个矩阵都行等价于一个
而且仅有一个简化阶梯形矩阵来结束本章。
这里的思想将在下一章中再次出现，
并得到进一步的发展。
% Here is an informal argument that the reduced 
% echelon form version of a matrix is unique.
% Consider again the example that started this section of a matrix that
% reduces to three different echelon form matrices.
% The first matrix of the three is the natural echelon form version.
% The second matrix is the same as 
% the first except that a row has been halved.
% The third matrix, too, is just a cosmetic variant of the first. 
% The definition of reduced echelon form outlaws this kind of fooling around.
% In reduced echelon form,
% halving a row is not possible because that would
% change the row's leading entry away from one, and
% neither is combining rows possible, because then a leading entry would no
% longer be alone in its column.

% This informal justification is not a proof;
% the argument shows that no two different reduced echelon form matrices
% are related by a single row operation step, but the argument does not
% rule out the possibility that two different reduced echelon form
% matrices could be related by multiple steps.
% Before we go to the proof, we finish this subsection by 
% rephrasing our work in a terminology that will be enlightening.
关键的观察涉及
行变换如何把一个矩阵变成另一个矩阵：
新的行是
原来的行的线性组合。

\begin{example}
考察这个高斯–若尔当消元。
\begin{align*}
  \begin{amat}{2}
    2  &1  &0  \\
    1  &3  &5
  \end{amat}
  &\grstep{-(1/2)\rho_1+\rho_2}
  \begin{amat}{2}
    2  &1    &0  \\
    0  &5/2  &5
  \end{amat}                           \\
  &\grstep[(2/5)\rho_2]{(1/2)\rho_1}
  \begin{amat}{2}
    1  &1/2  &0  \\
    0  &1    &2
  \end{amat}                        \\
  &\grstep{-(1/2)\rho_2+\rho_1}\;
  \begin{amat}{2}
    1  &0    &-1  \\
    0  &1    &2
  \end{amat}                       
\end{align*}
将这些矩阵记为 $A\rightarrow D\rightarrow G\rightarrow B$，
并把 $A$ 的行写作 $\alpha_1$ 与 $\alpha_2$，等等，我们就有下式。
\begin{align*}  % multline looks worse
  \left(\begin{array}{l}
    \alpha_1 \\
    \alpha_2
  \end{array}\right)
  &\grstep{-(1/2)\rho_1+\rho_2}
  \left(\begin{array}{l}
    \delta_1=\alpha_1 \\
    \delta_2=-(1/2)\alpha_1+\alpha_2
  \end{array}\right)                          \\
  &\grstep[(2/5)\rho_2]{(1/2)\rho_1}
  \left(\begin{array}{l}
    \gamma_1=(1/2)\alpha_1 \\
    \gamma_2=-(1/5)\alpha_1+(2/5)\alpha_2
  \end{array}\right)                          \\
  &\grstep{-(1/2)\rho_2+\rho_1}
  \left(\begin{array}{l}
    \beta_1=(3/5)\alpha_1-(1/5)\alpha_2  \\
    \beta_2=-(1/5)\alpha_1+(2/5)\alpha_2
  \end{array}\right)                         
\end{align*} 
\end{example}

\begin{example}
高斯操作对行作线性组合这一事实，
在含有行对换时同样成立。
对于这里的 \( A \)、\( D \)、\( G \) 与 \( B \)
\begin{equation*}
    \begin{mat}[r]
       0  &2  \\
       1  &1
     \end{mat}
    \grstep{\rho_1\leftrightarrow\rho_2}
    \begin{mat}[r]
       1  &1  \\
       0  &2
     \end{mat}                  
    \grstep{(1/2)\rho_2}
    \begin{mat}[r]
       1  &1  \\
       0  &1
     \end{mat}                   
    \grstep{-\rho_2+\rho_1}
    \begin{mat}[r]
       1  &0  \\
       0  &1
     \end{mat}
\end{equation*}
我们得到这些线性关系。
\begin{multline*}
  \left(\begin{array}{l}
    \vec{\alpha}_1 \\
    \vec{\alpha}_2
  \end{array}\right)
  \,\grstep{\rho_1\leftrightarrow\rho_2}\,
  \left(\begin{array}{l}
     \vec{\delta}_1 =\vec{\alpha}_2   \\
     \vec{\delta}_2 =\vec{\alpha}_1
  \end{array}\right)                                                 
  \,\grstep{(1/2)\rho_2}\,
  \left(\begin{array}{l}
     \vec{\gamma}_1 =\vec{\alpha}_2  \\
     \vec{\gamma}_2 =(1/2)\vec{\alpha}_1
  \end{array}\right)                                               \\
  \,\grstep{-\rho_2+\rho_1}\,
  \left(\begin{array}{l}
     \vec{\beta}_1 =(-1/2)\vec{\alpha}_1+1\cdot\vec{\alpha}_2   \\
     \vec{\beta}_2 =(1/2)\vec{\alpha}_1
  \end{array}\right)
\end{multline*}
\end{example}

总之，
高斯消元法系统地找出行的一个合适的
线性组合序列。

% \begin{definition}
% A \definend{linear combination}\index{linear combination} of 
% \( x_1,\ldots,x_m \)
% is an expression of the form $c_1x_1+c_2x_2+\,\cdots\,+c_mx_m$
% where the \( c \)'s are scalars.
%\end{definition}

% \noindent (We have already used the phrase 
% `linear combination' in this book.
% The meaning is unchanged, but the next result's statement makes
% a more formal definition in order.)

\begin{lemma}[线性组合引理] \label{lm:LinearCombinationLemma} \index{Linear Combination Lemma}
%<*lm:LinearCombinationLemma>
线性组合的线性组合仍是线性组合。
%</lm:LinearCombinationLemma>
\end{lemma}

\begin{proof}
%<*pf:LinearCombinationLemma>
给定由 $x$ 们的线性组合组成的集合，
从 $c_{1,1}x_1+\dots+c_{1,n}x_n$ 到 $c_{m,1}x_1+\dots+c_{m,n}x_n$，
考虑这些组合的一个组合
\begin{equation*}
  d_1(c_{1,1}x_1+\dots+c_{1,n}x_n)\,+\dots+\,d_m(c_{m,1}x_1+\dots+c_{m,n}x_n)
\end{equation*}
其中 $d$ 们与 $c$ 们一样都是标量。
将这些 $d$ 分配开来并重新组合，得到
\begin{equation*}
  %&=d_1c_{1,1}x_1+\dots+d_1c_{1,n}x_n\,
  % +d_2c_{2,1}x_1+\dots+\,
  % d_mc_{1,1}x_1+\dots+d_mc_{1,n}x_n         \\
  =(d_1c_{1,1}+\dots+d_mc_{m,1})x_1\,+\dots+\,(d_1c_{1,n}+\dots+d_mc_{m,n})x_n
\end{equation*}
它仍然是 $x$ 们的一个线性组合。
%</pf:LinearCombinationLemma>
\end{proof}


\begin{corollary} \label{cor:RowsOfEqMatsLinCombos}
%<*co:RowsOfEqMatsLinCombos>
当一个矩阵化简为另一个矩阵时，第二个矩阵的每一行
都是第一个矩阵的各行的线性组合。
%</co:RowsOfEqMatsLinCombos>
\end{corollary}

% The proof 
% uses induction. % \appendrefs{mathematical induction}\spacefactor=1000 %
% Before we proceed, here is an outline of the argument.
% For the base step, we
% will verify that the proposition is true when reduction 
% can be done in zero row operations.
% For the inductive step, we will 
% argue that if being able to reduce the first matrix to the second in some
% number $t\geq 0$ of operations implies that each row of the second is a linear
% combination of the rows of the first, then being able to reduce the first to
% the second in $t+1$ operations implies the same thing.
% Together these prove the result because  
% the base step shows that it is true in the zero operations case,
% and then the inductive step
% implies that it is true in the one operation case, and then the inductive step
% applied again gives that it is therefore true for two operations, etc.

\begin{proof}
%<*pf:RowsOfEqMatsLinCombos0>
对于任何两个可以互相化简的矩阵 $A$ 与~$B$，存在某个
把其中一个变为另一个的最少行变换次数。
我们对该次数进行归纳。

在基础步骤中，即可以用零次化简操作从一个矩阵变到另一个矩阵时，
这两个矩阵相等。
于是 $B$ 的每一行显然是 $A$ 的各行的组合：
$\vec{\beta}_i
  =0\cdot\vec{\alpha}_1+\cdots+1\cdot\vec{\alpha}_i+\cdots+0\cdot\vec{\alpha}_m$。
%</pf:RowsOfEqMatsLinCombos0>

%<*pf:RowsOfEqMatsLinCombos1>
归纳步骤假设归纳假设：~即当 $k\geq 0$ 时，
任何能由 \( A \) 经 \( k \) 次或更少操作得到的矩阵，
其各行都是 $A$ 的各行的线性组合。
考虑矩阵~$B$，使得把 \( A \) 化简为~\( B \)
需要$k+1$次操作。
在该化简中有一个倒数第二个矩阵~$G$，
即 $A\longrightarrow\cdots\longrightarrow G\longrightarrow B$。
归纳假设适用于这个 \( G \)，
因为它距 \( A \) 只有$k$步。
也就是说，\( G \) 的每一行
都是 \( A \) 的各行的线性组合。
%</pf:RowsOfEqMatsLinCombos1>

%<*pf:RowsOfEqMatsLinCombos2>
我们将验证~$B$ 的各行是~$G$ 的
各行的线性组合。
于是，线性组合引理（\nearbylemma{lm:LinearCombinationLemma}）
适用，它表明~$B$ 的各行是~$A$ 的
各行的线性组合。

如果把 \( G \) 变为~\( B \) 的行变换是对换，那么
$B$ 的行不过是 $G$ 的行的重新排列，$B$ 的每一行
都是 $G$ 的各行的线性组合。
如果把 $G$ 变为~$B$ 的操作是用标量~$c\rho_i$
乘某一行，
那么 $\vec{\beta}_i=c\vec{\gamma}_i$，其余各行不变。
最后，如果行变换是把一行的倍数加到
另一行上 $r\rho_i+\rho_j$，
那么 $B$ 中只有第~$j$行与~$G$ 的对应行不同，且
$\vec{\beta}_j=r\gamma_i+\gamma_j$，
它确实是 $G$ 的各行的线性组合。
%</pf:RowsOfEqMatsLinCombos2>

由于我们既证明了基础步骤又证明了归纳步骤，
根据数学归纳法原理，该命题得证。
\end{proof}

现在我们认识到，高斯消元法构造的是
行的线性组合。
但它的目标当然是终止于阶梯形，因为那是
线性方程组的一种特别基本的形态，
它把变量分离了出来。
例如，在这个矩阵
\begin{equation*}
  R=\begin{mat}[r]
    2  &3  &7  &8  &0  &0  \\
    0  &0  &1  &5  &1  &1  \\
    0  &0  &0  &3  &3  &0  \\
    0  &0  &0  &0  &2  &1
  \end{mat}
\end{equation*}
中，$x_1$ 已从 $x_5$ 的方程中被消去。
也就是说，高斯消元法使 $x_5$ 的行在某种意义上
独立于 $x_1$ 的行。

下面的结果使这个直觉变得精确。
我们有时把高斯消元法称作高斯消去（Gaussian elimination）。
它所消去的是各行之间的线性关系。

\begin{lemma}      \label{le:EchFormNoLinCombo}
%<*lm:EchFormNoLinCombo>
在阶梯形矩阵中，
没有哪个非零行是其余非零行的线性组合。
%</lm:EchFormNoLinCombo>
\end{lemma}

\begin{proof}
%<*pf:EchFormNoLinCombo0>
设 $R$ 是一个阶梯形矩阵，考察它的非 \( \vec{0} \) 行。
首先注意到，
如果把一行写成
其余各行的组合
$\vec{\rho}_i=c_1\vec{\rho}_1+\cdots+c_{i-1}\vec{\rho}_{i-1}+
               c_{i+1}\vec{\rho}_{i+1}+\cdots+c_m\vec{\rho}_m$
那么可以把该方程改写为
\begin{equation*}
   \vec{0}=c_1\vec{\rho}_1+\cdots+c_{i-1}\vec{\rho}_{i-1}+c_i\vec{\rho}_i+
               c_{i+1}\vec{\rho}_{i+1}+\cdots+c_m\vec{\rho}_m
  \tag{$*$}
\end{equation*}
其中并非所有系数都为零；具体地说，$c_i=-1$。
反过来也成立：~给定方程~($*$)，若某个 $c_i\neq 0$，则通过把
$c_i\vec{\rho}_i$ 移到左边并除以 $-c_i$，
我们可以把 $\vec{\rho}_i$ 表示成其余各行的组合。
因此，如果我们证明在~($*$) 中所有系数都是~$0$，
就证明了该定理。
为此我们对行号~$i$ 使用归纳法。
%</pf:EchFormNoLinCombo0>

%<*pf:EchFormNoLinCombo1>
基础情形是第一行~$i=1$
（如果没有这样的非零行，从而 $R$ 是零矩阵，则
引理平凡地成立）。
设 $\ell_i$ 是第~$i$行首主元所在的列号。
考察每一行位于第~$\ell_1$列的元素。
方程~($*$) 给出
\begin{equation*}
  0=c_1r_{1,\ell_1}+c_2r_{2,\ell_1}+\cdots+c_mr_{m,\ell_1}
  \tag{$**$}
\end{equation*}
该矩阵是阶梯形的，所以
第一行之后的每一行在该列的元素都是零：
$r_{2,\ell_1}=\cdots=r_{m,\ell_1}=0$。
于是方程~($**$) 表明 $c_1=0$，因为 $r_{1,\ell_1}\neq 0$，
它是该行的首主元。
%</pf:EchFormNoLinCombo1>

%<*pf:EchFormNoLinCombo2>
归纳步骤与基础步骤基本相同。
再次考察方程~($*$)。
我们将证明：如果对每个行标 $i\in\set{1,\ldots,k}$
系数 $c_i$ 都是 $0$，
那么 $c_{k+1}$ 也是 $0$。
我们关注来自第~$\ell_{k+1}$列的元素。
\begin{equation*}
  0=c_1r_{1,\ell_{k+1}}+\cdots+c_{k+1}r_{k+1,\ell_{k+1}}+\cdots+c_mr_{m,\ell_{k+1}}
\end{equation*}
由归纳假设，$c_1$、\ldots、$c_k$ 全是
$0$，于是上式化归为方程
$0=c_{k+1}r_{k+1,\ell_{k+1}}+\cdots+c_mr_{m,\ell_{k+1}}$。
该矩阵是阶梯形的，所以元素
$r_{k+2,\ell_{k+1}}$、\ldots、$r_{m,\ell_{k+1}}$ 全是~$0$。
于是 $c_{k+1}=0$，因为 $r_{k+1,\ell_{k+1}}\neq 0$，它是首主元。
%</pf:EchFormNoLinCombo2>
\end{proof}

有了这些，我们就可以证明高斯–若尔当消元的
最终产物是唯一的。

\begin{theorem}
\label{th:ReducedEchelonFormIsUnique}
%<*th:ReducedEchelonFormIsUnique>
每个矩阵都行等价于唯一的简化阶梯形矩阵。
%</th:ReducedEchelonFormIsUnique>
\end{theorem}

\begin{proof} \cite{Yuster}
%<*pf:ReducedEchelonFormIsUnique0>
固定行数 \( m \)。
我们将对列数 \( n \) 进行归纳。

基础情形是矩阵只有 \( n=1 \) 列。
如果它是零矩阵，那么它的阶梯形是零矩阵。
如果它含有非零元素，那么当该矩阵被化为
简化阶梯形时，它必含至少一个非零元素，而这个元素
必是第一行中的一个 \( 1 \)。
无论哪种情形，它的简化阶梯形都是唯一的。
%</pf:ReducedEchelonFormIsUnique0>

%<*pf:ReducedEchelonFormIsUnique1>
归纳步骤假设 \( n>1 \)，并且所有列数少于~\( n \) 的
\( m \)~行矩阵都有唯一的简化阶梯形。
考虑一个 \( \nbym{m}{n} \) 矩阵 \( A \)，并设
\( B \) 与 \( C \) 是由 \( A \) 得到的两个简化阶梯形矩阵。
我们将证明这两个矩阵必相等。
%</pf:ReducedEchelonFormIsUnique1>

%<*pf:ReducedEchelonFormIsUnique2>
设 \( \hat{A} \) 是由 \( A \) 的前 \( n-1 \) 列
组成的矩阵。
注意到
% if an $n$~column matrix is in reduced echelon form then any initial
% set of its columns is also in reduced echelon form, so \( \hat{A} \)
% is in reduced echelon form. 
任何把 \( A \) 化为简化
阶梯形的行变换序列也会把 \( \hat{A} \) 化为简化阶梯形。
由归纳假设，\( \hat{A} \) 的这个简化阶梯形
是唯一的，所以如果 \( B \) 与 \( C \) 不同，那么差异必
出现在第~\( n \) 列。
%</pf:ReducedEchelonFormIsUnique2>

我们通过证明这两者不可能只在该列不同，
来完成归纳步骤和整个论证。
%<*pf:ReducedEchelonFormIsUnique3>
考虑以 \( A \) 为系数矩阵的一个齐次方程组。
\begin{equation*}
  \begin{linsys}{4}
    a_{1,1}x_1  &+  &a_{1,2}x_2  &+  &\cdots  &+  &a_{1,n}x_n  &=  &0  \\
    a_{2,1}x_1  &+  &a_{2,2}x_2  &+  &\cdots  &+  &a_{2,n}x_n  &=  &0  \\
              &&&&&&&\vdotswithin{=}  \\
    a_{m,1}x_1  &+  &a_{m,2}x_2  &+  &\cdots  &+  &a_{m,n}x_n  &=  &0  
  \end{linsys}
  \tag{$*$}
\end{equation*}
根据定理~One.I.\ref{th:GaussMethod}
%the first theorem of this chapter 
该方程组的解集与 $B$ 的方程组
\begin{equation*}
  \begin{linsys}{4}
    b_{1,1}x_1  &+  &b_{1,2}x_2  &+  &\cdots  &+  &b_{1,n}x_n  &=  &0  \\
    b_{2,1}x_1  &+  &b_{2,2}x_2  &+  &\cdots  &+  &b_{2,n}x_n  &=  &0  \\
               &&&&&&&\vdotswithin{=}  \\
    b_{m,1}x_1  &+  &b_{m,2}x_2  &+  &\cdots  &+  &b_{m,n}x_n  &=  &0  
  \end{linsys}
  \tag{$**$}
\end{equation*}
的解集相同，也与 $C$ 的
\begin{equation*}
  \quad
  \begin{linsys}{4}
    c_{1,1}x_1  &+  &c_{1,2}x_2  &+  &\cdots  &+  &c_{1,n}x_n  &=  &0  \\
    c_{2,1}x_1  &+  &c_{2,2}x_2  &+  &\cdots  &+  &c_{2,n}x_n  &=  &0  \\
               &&&&&&&\vdotswithin{=}  \\
    c_{m,1}x_1  &+  &c_{m,2}x_2  &+  &\cdots  &+  &c_{m,n}x_n  &=  &0  
  \end{linsys}
  \tag{$\mathord{*}\mathord{*}\mathord{*}$}
\end{equation*}
%</pf:ReducedEchelonFormIsUnique3>
%<*pf:ReducedEchelonFormIsUnique4>
设 \( B \) 与 \( C \) 只在第~\( n \) 列不同，并假设
它们在第~\( i \) 行不同。
用 ($**$) 的第~\( i \) 行减去 ($\mathord{*}\mathord{*}\mathord{*}$) 的
第~\( i \) 行，
得到方程
\( (b_{i,n}-c_{i,n})\cdot x_n=0 \)。
我们已假设 \( b_{i,n}\neq c_{i,n} \)，于是得到 $x_n=0$。
因此 $x_n$ 不是自由变量，
于是在 ($**$) 和~($\mathord{*}\mathord{*}\mathord{*}$) 中
第 \( n \) 列含有某行的首主元，
因为在阶梯形矩阵中任何
不含首主元的列都对应于
一个自由变量。

但现在，由于 \( B \) 与 \( C \) 在
前 \( n-1 \)~列相等，根据
简化阶梯形的定义，它们在
第 \( n \) 列的首主元位于同一行。
而且，这两个首主元都必须是 $1$，
并且必须是该列中仅有的非零元素。
因此 \( B=C \)。
%</pf:ReducedEchelonFormIsUnique4>
\end{proof}

我们曾问过：
一个线性方程组的任意两个阶梯形版本
是否有相同数目的自由变量，如果有，
它们是否恰好是同样的变量？
有了前面的结果，我们可以对这两个问题都回答“是。”
不存在这样的线性方程组，使得
比如我们用一种方式应用高斯消元法得到 $y$ 与 $z$
是自由变量，而用另一种方式应用却得到 $y$ 与 $w$ 是自由变量。

在证明之前，
回顾一下自由变量与参数的区别。
这个方程组
\begin{equation*}
  \begin{linsys}{3}
    x &+ &y &  &  &= &1 \\ 
      &  &y &+ &z &= &2  
  \end{linsys}
\end{equation*}
有一个自由变量~$z$，因为它是唯一不引领一行的变量。
我们习惯用自由变量来参数化：$y=2-z$，$x=-1+z$，
但我们也可以用另一个变量参数化，例如
$z=2-y$，$x=1-y$。
所以参数的集合并不唯一，唯一的
是自由变量的集合。

\begin{corollary}
如果从一个初始线性方程组出发，我们用高斯消元法得到两个
不同的阶梯形方程组，那么这两个方程组有相同的自由变量。
\end{corollary}

\begin{proof}
前面的结果说简化阶梯形是唯一的。
从任一阶梯形版本出发，
向上消元即可得到简化阶梯形，
所以方程组的任一阶梯形版本与简化阶梯形版本
有相同的自由变量。
\end{proof}

我们以一段回顾来收尾。
在高斯消元法中，我们从一个矩阵出发，
然后导出一系列其他的矩阵。
如果我们能从一个矩阵导出另一个，我们就称这两个矩阵是相关的。
这种关系是一种等价关系，%\appendrefs{equivalence relation} 
称为行等价，
于是它把全体矩阵的集合划分为一些行等价类。
\begin{center}
  \includegraphics{gr/mp/ch1.30}
\end{center}
（图中所示的类中有无穷多个矩阵，但限于篇幅
我们只展示两个。）
我们已经证明，每个行等价类中
有且仅有一个简化阶梯形矩阵。
%<*ReducedEchelonFormIsCanonicalForm>
所以简化阶梯形是
行等价的一种规范形\appendrefs{canonical representatives}\spacefactor=1000
\index{canonical form!for row equivalence}\index{representative!for row equivalence classes}
：
简化阶梯形矩阵正是
这些类的代表。
%</ReducedEchelonFormIsCanonicalForm>
\begin{center}
  \includegraphics{gr/mp/ch1.31}
\end{center}

这里的思想是：理解一个数学情境的一种方式，
是能够对可能发生的各种情形加以分类。
这是本书的一个主题，
我们已经见过它好几次了。
我们曾把线性方程组的解集分为无元素、一个元素
与无穷多元素三种情形。
我们也曾把方程个数与未知量个数相同的线性方程组
分为非奇异与奇异两种情形。

这里，在我们考察行等价时，我们知道全体矩阵的
集合分裂成一些行等价类，
而且我们现在有办法确切指出每一个这样的类\Dash
我们可以把一个类中的矩阵看成是由该类中
唯一的简化阶梯形矩阵经行变换导出的。

用更操作性的说法，简化阶梯形的唯一性
使我们能够通过把关于类的问题
转化为关于代表的问题
来回答它们。
例如，正如本节开头所允诺的，现在我们可以
判定一个矩阵能否由另一个矩阵经行化简导出。
我们对两者都应用高斯–若尔当程序，看它们
是否给出相同的简化阶梯形。

\begin{example}  \label{ex:MatsNotRowEq}
这两个矩阵不是行等价的
\begin{equation*}
  \begin{mat}[r]
    1  &-3  \\
   -2  &6
  \end{mat}
  \qquad
  \begin{mat}[r]
    1  &-3  \\
   -2  &5
  \end{mat}
\end{equation*}
因为它们的简化阶梯形不相等。
\begin{equation*}
  \begin{mat}[r]
    1  &-3  \\
    0  &0
  \end{mat}
  \qquad
  \begin{mat}[r]
    1  &0   \\
    0  &1
  \end{mat}
\end{equation*}
\end{example}

\begin{example}
任何非奇异的 \( \nbyn{3} \) 矩阵经高斯–若尔当消元都化为
下面这个矩阵。
\begin{equation*}
    \begin{mat}[r]
      1  &0  &0 \\
      0  &1  &0 \\
      0  &0  &1
    \end{mat}
\end{equation*}
\end{example}

\begin{example} \label{ex:RowEqClassTwoTwoMats}
我们可以通过列出所有可能的
简化阶梯形矩阵来描述全部的类。
任何 $\nbyn{2}$ 矩阵都属于下面这些之一：~与这个矩阵
行等价的矩阵类，
\begin{equation*}
  \begin{mat}[r]
     0  &0  \\
     0  &0
  \end{mat}
\end{equation*}
与这种类型的某个矩阵行等价的矩阵构成的
无穷多个类
\begin{equation*}
  \begin{mat}
     1  &a  \\
     0  &0
  \end{mat}
\end{equation*}
其中 \( a\in\Re \)（包括 $a=0$），
与这个矩阵行等价的矩阵类，
\begin{equation*}
  \begin{mat}[r]
     0  &1  \\
     0  &0
  \end{mat}
\end{equation*}
以及与这个矩阵行等价的矩阵类
\begin{equation*}
  \begin{mat}[r]
     1  &0  \\
     0  &1
  \end{mat}
\end{equation*}
（这就是非奇异的 $\nbyn{2}$ 矩阵构成的类）。
\end{example}



\begin{exercises}
  \recommended \item 
    判断下列矩阵是否行等价。
    \begin{exparts*}
       \partsitem \(
           \begin{mat}[r]
             1  &2  \\
             4  &8
           \end{mat}, 
           \begin{mat}[r]
             0  &1  \\
             1  &2
           \end{mat} \)
       \partsitem \(
           \begin{mat}[r]
             1  &0  &2  \\
             3  &-1 &1  \\
             5  &-1 &5
           \end{mat},  
           \begin{mat}[r]
             1  &0  &2  \\
             0  &2  &10 \\
             2  &0  &4
           \end{mat} \)
       \partsitem \(
           \begin{mat}[r]
             2  &1  &-1 \\
             1  &1  &0  \\
             4  &3  &-1
           \end{mat},  
           \begin{mat}[r]
             1  &0  &2  \\
             0  &2  &10 \\
           \end{mat} \)
       \partsitem \(
           \begin{mat}[r]
             1  &1  &1  \\
            -1  &2  &2
           \end{mat},  
           \begin{mat}[r]
             0  &3  &-1 \\
             2  &2  &5
           \end{mat} \)
       \partsitem \(
           \begin{mat}[r]
             1  &1  &1  \\
             0  &0  &3
           \end{mat},  
           \begin{mat}[r]
             0  &1  &2  \\
             1  &-1 &1
           \end{mat} \)
    \end{exparts*}
    \begin{answer}
      把每个矩阵都化成简化阶梯形，然后进行比较。
      \begin{exparts}
        \partsitem 第一个给出
          \begin{equation*}
            \grstep{-4\rho_1+\rho_2}
            \begin{mat}[r]
              1  &2  \\
              0  &0
            \end{mat}
          \end{equation*}
          而第二个给出
          \begin{equation*}
            \grstep{\rho_1\leftrightarrow\rho_2}
            \begin{mat}[r]
              1  &2  \\
              0  &1
            \end{mat}
            \grstep{-2\rho_2+\rho_1}
            \begin{mat}[r]
              1  &0  \\
              0  &1
            \end{mat}
          \end{equation*}
          两个简化阶梯形矩阵并不相同，因此
          原来的矩阵不行等价。
        \partsitem 第一个如下。
          \begin{equation*}
            \grstep[-5\rho_1+\rho_3]{-3\rho_1+\rho_2}
            \begin{mat}[r]
              1  &0  &2  \\
              0  &-1 &-5 \\
              0  &-1 &-5
            \end{mat}
            \grstep{-\rho_2+\rho_3}
            \begin{mat}[r]
              1  &0  &2  \\
              0  &-1 &-5 \\
              0  &0  &0
            \end{mat}
            \grstep{-\rho_2}
            \begin{mat}[r]
              1  &0  &2  \\
              0  &1  &5  \\
              0  &0  &0
            \end{mat}
          \end{equation*}
          第二个如下。
          \begin{equation*}
            \grstep{-2\rho_1+\rho_3}
            \begin{mat}[r]
              1  &0  &2  \\
              0  &2  &10 \\
              0  &0  &0
            \end{mat}
            \grstep{(1/2)\rho_2}
            \begin{mat}[r]
              1  &0  &2  \\
              0  &1  &5  \\
              0  &0  &0
            \end{mat}
          \end{equation*}
          这两个矩阵行等价。
        \partsitem 这两个矩阵不行等价，因为它们的
          尺寸不同。
        \partsitem 第一个，
          \begin{equation*}
            \grstep{\rho_1+\rho_2}
            \begin{mat}[r]
              1  &1  &1  \\
              0  &3  &3
            \end{mat}
            \grstep{(1/3)\rho_2}
            \begin{mat}[r]
              1  &1  &1  \\
              0  &1  &1
            \end{mat}
            \grstep{-\rho_2+\rho_1}
            \begin{mat}[r]
              1  &0  &0  \\
              0  &1  &1
            \end{mat}
          \end{equation*}
          而第二个，
          \begin{equation*}
            \grstep{\rho_1\leftrightarrow\rho_2}
            \begin{mat}[r]
              2  &2  &5  \\
              0  &3  &-1
            \end{mat}
            \grstep[(1/3)\rho_2]{(1/2)\rho_1}
            \begin{mat}[r]
              1  &1  &5/2 \\
              0  &1  &-1/3
            \end{mat}
            \grstep{-\rho_2+\rho_1}
            \begin{mat}[r]
              1  &0  &17/6 \\
              0  &1  &-1/3
            \end{mat}
          \end{equation*}
          这两个矩阵不行等价。
        \partsitem 这里第一个是
          \begin{equation*}
            \grstep{(1/3)\rho_2}
            \begin{mat}[r]
              1  &1  &1  \\
              0  &0  &1
            \end{mat}
            \grstep{-\rho_2+\rho_1}
            \begin{mat}[r]
              1  &1  &0  \\
              0  &0  &1
            \end{mat}
          \end{equation*}
          而这是第二个。
          \begin{equation*}
            \grstep{\rho_1\leftrightarrow\rho_2}
            \begin{mat}[r]
              1  &-1 &1  \\
              0  &1  &2
            \end{mat}
            \grstep{\rho_2+\rho_1}
            \begin{mat}[r]
              1  &0  &3  \\
              0  &1  &2
            \end{mat}
          \end{equation*}
          这两个矩阵不行等价。
       \end{exparts}  
     \end{answer}
  \item 
    这些矩阵中哪些相互行等价？
    \begin{exparts*}
      \partsitem
         \(\begin{mat}
             1 &3 \\
             2 &4
           \end{mat}\)
      \partsitem
         \(\begin{mat}
             1 &5 \\
             2 &10
           \end{mat}\)
      \partsitem
         \(\begin{mat}
             1 &-1 \\
             3 &0
           \end{mat}\)
      \partsitem
         \(\begin{mat}
             2 &6 \\
             4 &10
           \end{mat}\)
      \partsitem
         \(\begin{mat}
             0  &1 \\
             -1 &0
           \end{mat}\)
      \partsitem
         \(\begin{mat}
             3 &3 \\
             2 &2
           \end{mat}\)
    \end{exparts*}
    \begin{answer}
      对每一个作高斯–若尔当消元。
      两个矩阵行等价当且仅当它们具有相同的
      简化阶梯形。
      下面是每个矩阵的简化形式。 
      % sage: M = matrix(QQ, [[1,3], [2,4]])
      % sage: gauss_jordan(M)
      % [1 3]
      % [2 4]
      %  take -2 times row 1 plus row 2
      % [ 1  3]
      % [ 0 -2]
      %  take -1/2 times row 2
      % [1 3]
      % [0 1]
      %  take -3 times row 2 plus row 1
      % [1 0]
      % [0 1]
      % sage: M = matrix(QQ, [[1,5], [2,10]])
      % sage: gauss_jordan(M)
      % [ 1  5]
      % [ 2 10]
      %  take -2 times row 1 plus row 2
      % [1 5]
      % [0 0]
      % sage: M = matrix(QQ, [[1,-1], [3,0]])
      % sage: gauss_jordan(M)
      % [ 1 -1]
      % [ 3  0]
      %  take -3 times row 1 plus row 2
      % [ 1 -1]
      % [ 0  3]
      %  take 1/3 times row 2
      % [ 1 -1]
      % [ 0  1]
      %  take 1 times row 2 plus row 1
      % [1 0]
      % [0 1]
      % sage: M = matrix(QQ, [[2,6], [4,10]])
      % sage: gauss_jordan(M)
      % [ 2  6]
      % [ 4 10]
      %  take -2 times row 1 plus row 2
      % [ 2  6]
      % [ 0 -2]
      %  take 1/2 times row 1
      %  take -1/2 times row 2
      % [1 3]
      % [0 1]
      %  take -3 times row 2 plus row 1
      % [1 0]
      % [0 1]
      % sage: M = matrix(QQ, [[0,1], [-1,0]])
      % sage: gauss_jordan(M)
      % [ 0  1]
      % [-1  0]
      %  swap row 1 with row 2
      % [-1  0]
      % [ 0  1]
      %  take -1 times row 1
      % [1 0]
      % [0 1]
      % sage: M = matrix(QQ, [[3,3], [2,2]])
      % sage: gauss_jordan(M)
      % [3 3]
      % [2 2]
      %  take -2/3 times row 1 plus row 2
      % [3 3]
      % [0 0]
      %  take 1/3 times row 1
      % [1 1]
      % [0 0]
      \begin{exparts*}
        \partsitem
          \(\begin{mat}
              1  &0  \\
              0  &1
            \end{mat} \)
        \partsitem
          \(\begin{mat}
              1  &5  \\
              0  &0
            \end{mat} \)
        \partsitem
          \(\begin{mat}
             1   &0  \\
             0   &1
            \end{mat} \)
        \partsitem
          \(\begin{mat}
             1   &0  \\
             0   &1
            \end{mat} \)
        \partsitem
          \(\begin{mat}
             1   &0  \\
             0   &1
            \end{mat} \)
        \partsitem
          \(\begin{mat}
             1   &1  \\
             0   &0
            \end{mat} \)
      \end{exparts*}
    \end{answer}
 \item 写出另外三个与所给矩阵行等价的矩阵。
   \begin{exparts*}
     \partsitem 
       \(\begin{mat}
           1 &3 \\
           4 &-1
         \end{mat}\)
     \partsitem
       \(\begin{mat}
           0 &1 &2 \\
           1 &1 &1  \\
           2 &3 &4
         \end{mat}\)
   \end{exparts*}
   \begin{answer}
     对每个小题，只需对起始矩阵作一些行
     变换即可。
      \begin{exparts}
        \partsitem
          把第一行乘以~$3$，得到下式。
          \begin{equation*}
            \begin{mat}
              3 &9  \\
              4 &-1
            \end{mat}
          \end{equation*}
          （这个行变换的选择并没有什么讲究；它 
          只是脑海中想到的第一个而已。）
          另外两个行变换是行对换 $\rho_1\leftrightarrow\rho_2$ 
          和倍加 $\rho_1+\rho_2$。
          \begin{equation*}
            \begin{mat}
              4 &-1  \\
              1 &3
            \end{mat}
            \quad
            \begin{mat}
              1 &3  \\
              5 &2
            \end{mat}
          \end{equation*}
        \partsitem 作同样的三个任意行变换，得到
          这三个。
          \begin{equation*}
           \begin{mat} 
            0 &3 &6 \\
            1 &1 &1  \\
            2 &3 &4  
           \end{mat}
           \quad         
           \begin{mat} 
            1 &1 &1  \\
            0 &1 &2 \\
            2 &3 &4  
           \end{mat}
           \quad         
           \begin{mat} 
            0 &1 &2 \\
            1 &2 &3  \\
            2 &3 &4  
           \end{mat}
           \quad         
          \end{equation*}
      \end{exparts}
    \end{answer}
  \recommended \item 对这个矩阵施行高斯消元法。
    把最终矩阵的每一行表示成起始矩阵
    各行的线性组合。
    \begin{equation*}
      \begin{mat}
        1 &2   &1 \\
        3 &-1  &0 \\
        0 &4   &0
      \end{mat}
    \end{equation*}
    \begin{answer}
      高斯化简是常规计算。
      % sage: load("gauss_method.sage")
      % sage: M = matrix (QQ, [[1,2,1], [3,-1,0], [0,4,0]])
      % sage: gauss_method(M)
      % [ 1  2  1]
      % [ 3 -1  0]
      % [ 0  4  0]
      %  take -3 times row 1 plus row 2
      % [ 1  2  1]
      % [ 0 -7 -3]
      % [ 0  4  0]
      %  take 4/7 times row 2 plus row 3
      % [    1     2     1]
      % [    0    -7    -3]
      % [    0     0 -12/7]
      \begin{equation*}
        \begin{mat}
          1 &2   &1 \\
          3 &-1  &0 \\
          0 &4   &0
        \end{mat}
        \grstep{-3\rho_1+\rho_2}    
        \begin{mat}
          1 &2   &1 \\
          0 &-7  &-3 \\
          0 &4   &0
        \end{mat}
        \grstep{(4/7)\rho_2+\rho_3}    
        \begin{mat}
          1 &2   &1 \\
          0 &-7  &-3 \\
          0 &0   &-12/7
        \end{mat}
      \end{equation*}
      把这三个矩阵分别记为 $A$、$D$ 和~$B$，则有下式。
      \begin{align*}
        \begin{mat}
          \alpha_1 \\
          \alpha_2 \\
          \alpha_3
        \end{mat}
        &\grstep{-3\rho_1+\rho_2}    
        \begin{mat}
          \delta_1=\alpha_1 \\
          \delta_2=-3\alpha_1+\alpha_2 \\
          \delta_3=\alpha_3
        \end{mat}                                   \\
        & \grstep{(4/7)\rho_2+\rho_3}
        \begin{mat}
          \beta_1=\alpha_1  \\
          \beta_2=-3\alpha_1+\alpha_2 \\
          \beta_3=-(12/7)\alpha_1+(4/7)\alpha_2+\alpha_3
        \end{mat}
      \end{align*}
    \end{answer}
  \item 
     描述
     \nearbyexample{ex:RowEqClassTwoTwoMats} 所展示的各个等价类中的矩阵。
     \begin{answer}
       首先，与全
       \( 0 \) 矩阵行等价的矩阵只有它自己（因为行变换对它没有效果）。

       其次，能化简为 
       \begin{equation*}
         \begin{mat}
           1  &a  \\
           0  &0
         \end{mat}
       \end{equation*}
       的矩阵形如
       \begin{equation*}
         \begin{mat}
           b  &ba \\
           c  &ca
         \end{mat}
       \end{equation*}
       （其中 \( a,b,c\in\Re \)，且 \(b\) 与 \(c\) 不全为零）。  

       再次，能化简为 
       \begin{equation*}
         \begin{mat}[r]
           0  &1  \\
           0  &0
         \end{mat}
       \end{equation*}
       的矩阵形如
       \begin{equation*}
         \begin{mat}
           0  &a \\
           0  &b
         \end{mat}
       \end{equation*}
       （其中 \( a,b\in\Re \)，且不全为零）。  

       最后，能化简为 
       \begin{equation*}
         \begin{mat}[r]
           1  &0  \\
           0  &1
         \end{mat}
       \end{equation*}
       的矩阵是非奇异的矩阵。
       这是因为以这个矩阵为系数矩阵的线性方程组将有唯一解，而这正是
       非奇异的定义。
       （换一种说法：它们不属于上面的任何一个等价类。）
     \end{answer}
  \item 
    描述与下列矩阵
    同一行等价类中的所有矩阵。
    \begin{exparts*}
       \partsitem  \(
           \begin{mat}[r]
             1  &0  \\
             0  &0
           \end{mat}  \)
       \partsitem  \(
           \begin{mat}[r]
             1  &2      \\
             2  &4
           \end{mat} \)
       \partsitem \(
           \begin{mat}[r]
             1  &1      \\
             1  &3
           \end{mat} \)
    \end{exparts*}
    \begin{answer}
      \begin{exparts}
        \partsitem 它们形如
          \begin{equation*}
            \begin{mat}
              a  &0  \\
              b  &0
            \end{mat}
          \end{equation*}
          其中 \( a,b\in\Re \) 至少有一个不为零。
        \partsitem 它们具有如下形式 
          \begin{equation*}
            \begin{mat}
             a  &2a \\
             b  &2b
            \end{mat}
          \end{equation*}
          其中 \( a,b\in\Re \) 至少有一个不为零。
        \partsitem 所给矩阵是非奇异的。
          所以该行等价类由所有非奇异的 $\nbyn{2}$
          矩阵组成。
          （给出公式的话，它们形如
          \begin{equation*}
            \begin{mat}
              a  &b  \\
              c  &d
            \end{mat}
          \end{equation*}
          其中 \( a,b,c,d\in\Re \)）且 \( ad-bc\neq 0 \)。
          我们将在第四章看到，这个公式 
          判定了 \( \nbyn{2} \) 矩阵
          何时非奇异。）
      \end{exparts}  
    \end{answer}
  \item 
    行等价类有多少个？
    \begin{answer}
       有无穷多个。
       例如，在 
       \begin{equation*}
         \begin{mat}
           1  &k  \\
           0  &0
         \end{mat}
       \end{equation*}
       中，每个 $k\in\Re$ 给出一个不同的等价类。  
    \end{answer}
  \item 
    行等价类能包含不同尺寸的矩阵吗？
    \begin{answer}
      不能。
      行变换不改变矩阵的尺寸。  
    \end{answer}
  \item 
    行等价类有多大？
    \begin{exparts} 
      \partsitem 证明：对任意全零矩阵，其等价类是有限的。
      \partsitem 还有别的等价类只含有限多个成员吗？
    \end{exparts}
    \begin{answer}
     \begin{exparts}
      \partsitem 对全零矩阵作行变换没有效果。
        因此每个这样的矩阵在它的行等价类中只有自己。  
      \partsitem 没有。
        任何非零元都可以倍乘。
     \end{exparts}
    \end{answer}
  \recommended \item 
    给出两个首主元位于相同列的
    简化阶梯形矩阵，
    但它们并不行等价。
    \begin{answer}
      这里有两个。
      \begin{equation*}
        \begin{mat}[r]
          1  &1  &0  \\
          0  &0  &1
        \end{mat}
        \quad\text{和}\quad
        \begin{mat}[r]
          1  &0  &0  \\
          0  &0  &1
        \end{mat}
      \end{equation*}  
     \end{answer}
  \recommended \item 
    证明任意两个 \( \nbyn{n} \) 非奇异矩阵
    是行等价的。
    任意两个奇异矩阵也行等价吗？
    \begin{answer}
      任意两个 \( \nbyn{n} \) 非奇异矩阵具有
      相同的简化阶梯形，即除对角线上
      为 \( 1 \) 之外全为 \( 0 \) 的那个矩阵。
      \begin{equation*}
        \begin{mat}
          1  &0  &       &0  \\
          0  &1  &       &0  \\
             &   &\ddots &   \\
          0  &0  &       &1
        \end{mat}
      \end{equation*}

      两个同尺寸的奇异矩阵未必行等价。
      例如，这两个 \( \nbyn{2} \) 奇异矩阵
      不行等价。
      \begin{equation*}
        \begin{mat}[r]
          1  &1  \\
          0  &0
        \end{mat}
        \quad\text{和}\quad
        \begin{mat}[r]
          1  &0  \\
          0  &0
        \end{mat}
      \end{equation*}  
    \end{answer}
  \recommended \item 
    描述包含下列矩阵的所有行等价类。
    \begin{exparts*}
      \partsitem \( \nbyn{2} \)~矩阵
      \partsitem \( \nbym{2}{3} \)~矩阵
      \partsitem \( \nbym{3}{2} \)~矩阵
      \partsitem \( \nbyn{3} \)~矩阵
    \end{exparts*}
    \begin{answer}
      由于每个等价类中有且仅有一个简化阶梯形矩阵，
      我们只需列出所有可能的简化阶梯形矩阵。

      这份清单参见 \nearbyexercise{exer:PossRedEchFrms} 的答案。 
    \end{answer}
  \item  
     \begin{exparts}
          \partsitem 证明：向量 $\vec{\beta}_0$ 是集合 $\set{\vec{\beta}_1,\ldots,\vec{\beta}_n}$
            中成员的线性组合当且仅当存在线性关系 
            $\zero=c_0\vec{\beta}_0+\cdots+c_n\vec{\beta}_n$
            其中 $c_0$ 不为零。
            （\textit{提示。}   当心 $\vec{\beta}_0=\zero$ 的情形。）
          \partsitem 利用它来简化 
             \nearbylemma{le:EchFormNoLinCombo} 的证明。   
       \end{exparts}
       \begin{answer}
          \begin{exparts}
           \partsitem 如果存在 $c_0$ 不为零的线性关系
             那么我们可以从两边减去 $c_0\vec{\beta}_0$ 再除以 $-c_0$，
             把 $\vec{\beta}_0$ 表示成其余向量的线性
             组合。
             （注：  
             如果集合中没有其他向量\Dash 即关系式
             是 $\zero=3\cdot\zero$\Dash 那么该命题仍然成立，因为
             按定义零向量是空向量集 
             的和。）

             反之，如果 $\vec{\beta}_0$ 是其余向量的组合 
             $\vec{\beta}_0=c_1\vec{\beta}_1+\dots+c_n\vec{\beta}_n$
             那么从两边减去
             $\vec{\beta}_0$ 就给出一个系数
             不全为零的线性关系；即
             $\vec{\beta}_0$ 前面的 $-1$。
           \partsitem 第一行不是其余各行的线性组合，
             理由
             如证明中所述：~在包含第一行首主元的那一列
             的分量等式中，唯一非零的元素
             是来自第一行的首主元，因此
             它的系数必为零。
             于是，由本题前面的小题，第一行与其
             余各行没有线性关系。

             因此，在考虑第二行能否与其余各行有线性 
             关系时，
             我们可以不考虑第一行。
             但这时刚才对第一行的论证也将适用于
             第二行。
             （也就是说，我们这里是用归纳法论证的。）             
         \end{exparts}
      \end{answer}
  % \item 
  %   Why, in the proof of \nearbytheorem{th:ReducedEchelonFormIsUnique},
  %   do we bother to restrict to the nonzero rows?
  %   Why not just stick to the relationship that we began with,
  %   $\beta_i=c_{i,1}\delta_1+\dots+c_{i,m}\delta_m$, with $m$ instead of $r$,
  %   and argue using it that the only nonzero coefficient
  %   is \( c_{i,i}  \), which is \( 1 \)?
  %   \begin{answer}
  %      The zero rows could have nonzero coefficients, and
  %      so the statement would not be true.
  %   \end{answer}
  \recommended \item 
   \cite{Trono}
   三位卡车司机走进一家路边咖啡馆。
   一位卡车司机买了四个三明治、一杯咖啡和十个 
   甜甜圈，花了 \$$8.45$。
   另一位司机买了三个三明治、一杯咖啡和七个
   甜甜圈，花了 \$$6.30$。
   第三位卡车司机买一个三明治、一杯咖啡和 
   一个甜甜圈要付多少钱？
    \begin{answer}
      我们知道 $4s+c+10d=8.45$ 与 $3s+c+7d=6.30$，而且想知道
      $s+c+d$ 是多少。
      幸运的是，$s+c+d$ 是 $4s+c+10d$ 与 $3s+c+7d$ 的线性组合。
      把未知的价格记为 $p$，我们有如下化简过程。
      \begin{align*}
        \begin{amat}{3}
          4  &1  &10  &8.45 \\
          3  &1  &7   &6.30 \\
          1  &1  &1   &p
        \end{amat}
        &\grstep[-(1/4)\rho_1+\rho_3]{-(3/4)\rho_1+\rho_2}
        \begin{amat}{3}
          4  &1    &10     &8.45      \\
          0  &1/4  &-1/2   &-0.037\,5 \\
          0  &3/4  &-3/2   &p-2.112\,5
        \end{amat}                                      \\
        &\grstep{-3\rho_2+\rho_3}
        \begin{amat}{3}
          4  &1    &10     &8.45      \\
          0  &1/4  &-1/2   &-0.037\,5 \\
          0  &0    &0      &p-2.00
        \end{amat}
     \end{align*}
     所付的价钱是 \$$2.00$。
   \end{answer}
   % \item
   %  The fact that Gaussian reduction disallows multiplication of
   %  a row by zero is needed for the proof of uniqueness of reduced echelon form,
   %  or else every matrix would
   %  be row equivalent to a matrix of all zeros.
   %  Where is it used?
   %  \begin{answer}
   %    If multiplication of a row by zero were allowed then
   %    \nearbylemma{le:EquivMatsSameForm}
   %    would not hold.
   %    That is, where
   %    \begin{equation*}
   %      \begin{mat}[r]
   %        1  &3  \\
   %        2  &1
   %      \end{mat}
   %      \grstep{0\rho_2}
   %      \begin{mat}[r]
   %        1  &3  \\
   %        0  &0
   %      \end{mat}
   %    \end{equation*}
   %    all the rows of the second matrix can be expressed as linear combinations
   %    of the rows of the first, but the converse does not hold.
   %    The second row of the first matrix is not a linear combination of the
   %    rows of the second matrix.  
   %  \end{answer}
  \item 
   线性组合引理说的是，对给定的线性方程组作
   高斯消元能得到哪些方程。
   \begin{enumerate}
     \item 写出一个不被这个方程组蕴含的方程。
       \begin{equation*}
         \begin{linsys}{2}
           3x  &+  &4y  &=  &8 \\
           2x  &+  & y  &=  &3 
         \end{linsys}
       \end{equation*}
     \item 从无解的方程组能导出任何方程吗？
   \end{enumerate}
   \begin{answer}
     \begin{enumerate}
        \item 一个容易的答案是这个。
          \begin{equation*}
            0=3
          \end{equation*}
          想要一个不那么抬杠的答案，求解该方程组：
          \begin{equation*}
            \begin{amat}[r]{2}
              3  &4  &8  \\
              2  &1   &3
            \end{amat}
            \grstep{-(2/3)\rho_1+\rho_2}
            \begin{amat}[r]{2}
              3  &4    &8    \\
              0  &-5/3 &-7/3
            \end{amat}
          \end{equation*}
          得到 \( y=7/5 \) 与 \( x=4/5 \)。
          于是任何不被 \( (7/5,4/5) \) 满足的方程都可以，
          例如 \( 5x+5y=10 \)。
        \item 从无解的方程组可以导出每一个方程。
          例如，下面演示如何从
          \( 0=5 \) 导出 \( 3x+2y=4 \)。
          首先，
          \begin{equation*}
            0=5
            \grstep{(3/5)\rho_1}
            0=3
            \grstep{x\rho_1}
            0=3x
          \end{equation*}
          （\( x=0 \) 情形的有效性需要另外论证，但很显然）。
          类似地，\( 0=2y \)。
          对 \( 0=4 \) 亦然。
          而现在，\( 0+0=0 \) 便给出 \( 3x+2y=4 \)。
     \end{enumerate}  
    \end{answer}
  \item 
    \cite{HoffmanKunze}
    把行等价的定义推广到线性方程组。
    在你的定义下，等价的方程组有相同的解集吗？
    \begin{answer}
      定义：若它们的增广
      矩阵行等价，则称线性方程组等价。
      等价方程组有相同解集的证明很容易。  
    \end{answer}
  \item 
    在这个矩阵
    \begin{equation*}
      \begin{mat}[r]
        1  &2  &3  \\
        3  &0  &3  \\
        1  &4  &5
      \end{mat}
    \end{equation*}
    中，第一列与第二列之和等于第三列。
    \begin{exparts}
      \partsitem 证明在任何行变换下这一点仍然成立。
      \partsitem 提出一个猜想。
      \partsitem 证明它成立。
    \end{exparts}
    \begin{answer}
      \begin{exparts}
        \partsitem 三种可能的行对换很容易， 
          三种可能的倍乘也是如此。
          六种可能的倍加之一是 \( k\rho_1+\rho_2 \)：
          \begin{equation*}
            \begin{mat}
              1           &2           &3  \\
              k\cdot 1+3  &k\cdot 2+0  &k\cdot 3+3  \\
              1           &4           &5
            \end{mat}
          \end{equation*}
          此时第一列与第二列之和仍等于第三列。
          其余五种倍加类似。
        \partsitem 显然的猜想是：行变换不改变
          列之间的线性关系。
        \partsitem 逐情形的
          证明按照第一小题给出的思路进行。
      \end{exparts}  
   \end{answer}
\end{exercises}
